\section{Part B. Limit actions and reduction}
\label{sec:B}

Unless a different hypothesis is stated, $R$ is a henselian discrete valuation
ring of residue characteristic zero, $K$ is its fraction field, and the residue
field $k$ is finitely generated over $\Q$. We distinguish
$E_{U_K}=\pi_1(U_K)$ from the fundamental group of a smooth model over $R$.
This distinction is essential in the contact calculation below.
The manuscript anchor for B0--B7 and B9--B11 is cote 149, manuscript
pp.~49--51, scans 99--104, especially equations (30)--(31) and conditions
a)--f). We use the author's page numbers at the top of recto pages; the
sequential numbers at the lower left are a different pagination.

\subsection{B0: comparison and the realization obstruction}

\begin{claim}{B.1}{Proved}
For a torus $S$ and an abelian variety $A$ over a field $k$ finitely generated
over $\Q$, $\Hom_{G_k}(TS,TA)=0$. Also
$\Hom_{G_k}(V_\ell A,\Q_\ell)=0$ for every prime $\ell$.
\end{claim}


\begin{proof}
We may replace $k$ by a finite extension splitting $S$: a map nonzero before
restriction is still nonzero afterwards. A nonzero continuous map of full
Tate modules has a nonzero $\ell$-component for some $\ell$; components with
different residue primes have no nonzero continuous homomorphisms. Thus it
would give a nonzero $G_k$-equivariant map $\Z_\ell(1)\to T_\ell A$.
Choose a smooth integral model $M$ of $k$ over its algebraic constant number
field, and shrink $M$ until $A$ extends to an abelian scheme. Both local
systems factor through $\pi_1(M)$, and $G_k\twoheadrightarrow\pi_1(M)$ is
surjective because a connected finite etale cover of normal $M$ has connected
generic fibre. Equivariance therefore holds for $\pi_1(M)$.
At a closed point $z$ of $M$, the pullback of the proposed map, after an
etale path identifies geometric fibres, is still nonzero and is
$G_{\kappa(z)}$-equivariant. The field $\kappa(z)$ is a number field.
This contradicts Stix's Lemma 19, which states the torus--abelian variety
vanishing over the number fields in question \citep[Lemma 19]{CF_Stix2008}.
The result invoked there is Ribet's Theorem 1: for an abelian variety over a
number field, its torsion over the extension obtained by adjoining all roots
of unity is finite \citep[Appendix, Theorem 1]{BC_KatzLang}.
Indeed a nonzero image of $\Z_\ell(1)$ would give infinitely many torsion
points over that extension. This verifies the precise implication from the
opened theorem rather than treating an earlier AI return as authority.

For the second assertion the Weil pairing identifies
$(V_\ell A)^*$ with $V_\ell A^\vee(-1)$. A map to $\Q_\ell$ is consequently
a map $\Q_\ell(1)\to V_\ell A^\vee$; multiplication by an appropriate power
of $\ell$ makes it integral. The first assertion excludes it.
\end{proof}

\begin{claim}{B.2}{Proved}
Let $1\to\Pi\to E\to G\to1$ be a continuous extension of profinite groups
with $\Pi$ topologically finitely generated and centre-free. Its outer action
$\rho:G\to\Out(\Pi)$ determines the extension: the map
\[
 E\longrightarrow\Aut(\Pi)\times_{\Out(\Pi)}G,
 \qquad e\longmapsto(\operatorname{conj}_e,\bar e)
\]
is an isomorphism of profinite extensions. This is the part of F149-C4 used
in R6 and R11; it requires no 2026 theorem about relative curves.
\end{claim}
\begin{proof}
The topology on $\Aut(\Pi)$ is the inverse-limit topology on its actions on
finite characteristic quotients. Finite generation makes those quotients
cofinal among finite quotients, so $\Aut(\Pi)$ is profinite. The displayed map
is continuous. Its kernel is $Z(\Pi)=1$. Given $(a,g)$ on the right, lift
$g$ to $e\in E$. Equality of outer classes gives
$a=\operatorname{int}(h)\operatorname{conj}_e$ for some $h\in\Pi$; hence
$(a,g)$ is the image of $he$. A continuous bijection between compact
Hausdorff groups is a topological isomorphism. In the applications the
geometric group is a hyperbolic surface group or a nonabelian free profinite
group. Centre-freeness in both cases, including every nonempty prime set,
is the explicit statement of \citet[Remark 1.1.3]{BC_MochComb}.
\end{proof}

\begin{claim}{B.3}{Known}
For a stable marked model in the stated characteristic-zero setting, its log
geometric group is $\Pi_{\mathcal G}\simeq\Pi_U$. A geometric node of thickness
$m$ gives an injected edge subgroup $\Pi_e\simeq\widehat\Z(1)$. After an
unramified extension defining the node and its branches, the node log group
and its map to the base log group are
\[
 T_P\rtimes G_k\longrightarrow T\rtimes G_k,
 \qquad (a,b,c,g)\longmapsto(c,g),\qquad
 T_P=\{(a,b,c):a+b=mc\}.
\]
Its geometric kernel is the antidiagonal $\{(a,-a,0)\}$, not the whole
rank-two node inertia group. Every residue splitting makes this antidiagonal
essentially cyclotomically normalized.
\end{claim}
\begin{proof}
Here are the comparison inputs, with their hypotheses. Illusie's Theorem 7.6
identifies finite Kummer etale covers of a log regular fs scheme with finite
etale covers of its trivial-log open tamely ramified at the boundary's
codimension-one points \citep[Theorem 7.6]{BC_Illusie}. The published proof is
Mochizuki's Log Purity Theorem 3.3: normalization in such a cover is log
regular and the map extends uniquely log etale
\citep[Theorem 3.3]{BC_MochExt}. Hoshi explicitly checks that the constructed
chart $P\subset P_Y\subset n^{-1}P$ makes this map Kummer
\citep[Remark 1.10]{BC_Hoshi}. Full faithfulness follows because maps between
normalizations are determined on their dense interiors and integral closure
extends every map of their function fields. Hoshi's Corollary 3.9, in the
numbering of RIMS-1558 (July 2006), identifies Kummer etale cover categories
of a proper log regular fs scheme over a complete local ring and its strict
closed log fibre. These statements apply to the curve with the divisorial
log structure of its special fibre and marks. In residue characteristic zero
all the relevant ramification is tame.

Completion does not alter the groups being compared. On absolute Galois
groups the natural restriction $G_{\widehat K}\to G_K$ is an isomorphism:
it is the identity on $G_k$, and on inertia it identifies the compatible
Kummer roots of a uniformizer, giving the identity on $T$. The Kummer
classification follows from SGA 1, XIII, Propositions 5.1--5.2: adjoining
prime-to-residue-characteristic roots of normal-crossing parameters removes
tame ramification. On geometric groups the base-extension map is an
isomorphism in characteristic zero. The exact comparison used here is
SGA 1, XIII, Proposition 4.3 (homotopy exactness for the stated locally
1-aspherical morphisms with a section), Examples 4.4 (smooth proper pairs
with relative normal-crossing boundary), and 4.5.3 (base-change compatibility
of the relative pro-$\Sigma$ group) \citep{BC_SGA1}. A geometric section is
available after passing to the algebraically closed base; the arithmetic
map then has isomorphic kernel and quotient. These maps transport the
specified splitting, not merely its existence.
The decisive wording in XIII, 4.5.3 is ``canoniquement isomorphe'' to the
pulled-back relative group; in XIII, 5.2 the extension is ``de mani\`ere
unique \`a isomorphisme unique pr\`es''. The regularity and normal-crossing
hypotheses of 5.1--5.2 are applied to the trait parameter; the possibly
nonregular node of thickness $m$ uses log purity and the saturated chart
below, not an illicit application of regular-scheme purity to that node.

The model is a stable log curve because its classifying morphism to the
stack of stable marked curves pulls back the universal log curve; at a node
the smoothing parameter maps to a unit times $\varpi^m$. Locally its chart is
\[
 P=\langle X,Y,T_0\mid X+Y=mT_0\rangle.
\]
This monoid is saturated: in $P^{\rm gp}=\Z^2$, with basis $(X,T_0)$, its
cone consists of the integral pairs $(r,s)$ with $s\ge0$ and $s+mr\ge0$;
if $r\ge0$ use $X,T_0$, and if $r<0$ use $Y=(-1,m),T_0$.
Kummer characters are therefore $\Hom(P^{\rm gp},T)=T_P$.
Compatible roots satisfy $x^{1/n}y^{1/n}=(\varpi^{1/n})^m$; this proves the
displayed homomorphism and its cyclotomic action coordinate by coordinate.

Mochizuki's Example 2.5 identifies the geometric admissible group with the
semi-graph group \citep[Example 2.5]{BC_MochComb}. The admissible-cover
description in \citet[\S2, Definition 2.1]{BC_Moch1996} uses covers which
become admissible after a tame base extension. In characteristic zero the
same local equations are the Kummer charts above: on a normalized component
a branch has the equation $z=w^n$, while the paired branch has opposite
character. Matching those branches and the common root of the smoothing
parameter is exactly both gluing conditions. Normalization and the
full-faithfulness already proved identify the resulting categories, including
their morphisms.

For the minimal deformation log point the node-group injection is
Hoshi--Mochizuki's Lemma 2.11(ii); part (iii) gives $(a,b)\mapsto a+b$
\citep[Lemma 2.11]{BC_HMNodes}. Pull back the extension along
$T\to\bigoplus_eT$, $c\mapsto(m_ec)_e$, dual to
$\bigoplus_e\mathbb N\to\mathbb N$, $(n_e)\mapsto\sum m_en_e$. An injection remains an
injection after a fibre product. Its kernel is the displayed antidiagonal
and maps onto the edge-like subgroup. This proves the assertion for arbitrary
thickness, without identifying the edge with the rank-two inertia.

Finally, the inverse image of $s(G_k)$ in $T_P\rtimes G_k$ is a closed
subgroup which maps onto $s(G_k)$: lift $c$ by $(mc,0,c)$. Conjugation on
its kernel is multiplication by $\chi(g)$, since $T_P$ is abelian. Taking
its image in the arithmetic curve group proves essential cyclotomic
normalization for every $s$. No splitting of that inverse image was assumed.
\end{proof}

\begin{claim}{B.4}{Proved}
R10 and R11 hold with the node notation $\Pi_e$: a stable nodal fibre
prevents $E_s$ from being the extension of a smooth hyperbolic curve after
any finite residue extension; for arbitrary $U$, realization of $\rho_s$
over $k$, realization after a finite extension of $k$, and potentially good
marked reduction are equivalent. Consequently B9 holds.
\end{claim}
\begin{proof}
Mochizuki's Proposition 1.2 makes distinct edge-like subgroups distinct up to
geometric conjugation and commensurably terminal. In particular a procyclic
overgroup normalizes an edge group and must equal it. Nodes and cusps thus
give $\nu+n$ distinct maximal procyclic classes when there are $n$ nodes.
By B.1, Stix's Theorem 20 applies to a smooth curve over $k$: nontrivial
essentially cyclotomically normalized procyclic subgroups are precisely
subgroups of cusp inertia. His Proposition 11 and Corollary 13 give their
normalizers and distinctness, and Corollary 21 gives precisely one maximal
class per geometric cusp \citep{CF_Stix2008}.

We must exclude a smooth realizing curve of a different type. Let
$t=b_1(\mathcal G)$. The graphic quotient of $H_1(\Pi,\Q_\ell)$ has dimension
$t$ and becomes trivial after a finite residue extension, for every splitting:
inertia acts trivially on the graph. This is also Proposition 2.4(i) of
\citet{BC_MochComb}. The homology of a smooth curve has an exact sequence
\[
 0\to\Q_\ell(1)^{\max(\nu-1,0)}\to
 H_1\to V_\ell J\to0
\]
after the cusps are split. B.1 and the infinite cyclotomic image show it has
no trivial quotient. Hence $t>0$ is impossible for a smooth realization.
For a tree, $H^2(\Pi,\Z/\ell)$ distinguishes proper and affine: it is
$\Z/\ell$ for the profinite genus-$g$ surface group, and zero for the free
profinite group of a punctured hyperbolic surface. Here is the required
comparison. In the tower of all finite etale covers of a smooth proper
hyperbolic curve over an algebraically closed characteristic-zero field,
every $H^1(-,\Z/\ell)$ class dies on the cover representing its torsor.
Every $H^2$ class dies after a further cover of degree $\ell$, since
pullback multiplies the orientation class by the degree. Such a cover
exists at every stage because the genus is at least two. Higher cohomology
vanishes. The direct limit in positive degree is therefore zero, so the
Cartan--Leray comparison identifies continuous group cohomology with the
curve's etale cohomology. Poincare duality gives $H^2=\Z/\ell$;
this is also the surface duality of \citet[Proposition 1.3]{BC_MochComb}.
For an affine hyperbolic curve SGA 1, XIII, Corollary 2.12 gives a free
profinite group. It is projective: lift its free generators through any
finite epimorphism. Hence its continuous $H^2$ with finite coefficients
vanishes. These are intrinsic group invariants.

In the proper tree case a single node already contradicts Stix's theorem.
In the affine tree case, eliminate leaf vertices in the boundary homology
relations: each node class becomes the sum of the cusp classes on one side,
leaving exactly the single relation among the $\nu$ cusp classes. Therefore
\[
 0\to\Q_\ell(1)^{\nu-1}\to H_1(\Pi,\Q_\ell)
 \to\bigoplus_v V_\ell\Jac(C_v)\to0.
\]
By B.1, the dimension of $\Hom_{G_k}(\Q_\ell(1),H_1)$ is $\nu-1$.
A smooth realizing affine curve has the same dimension, hence the same
number $\nu$ of cusps, contradicting the $\nu+n$ classes. This also proves
the counting criterion in L5, including $\nu=0$. A genuine point or cusp
action has the outer action of its smooth curve; forgetting the lift proves
B9 from this obstruction.

For completeness we supply the finite-extension comparison used in the
potential statement. For a finite extension $K_1/K$ inside the chosen
separable closure put
\[
 H=s^{-1}(G_{K_1}),\qquad G_L=G_{K_1}\cap p^{-1}(H),
 \qquad p:G_K\to G_k.
\]
Then $H$ is open, $s(H)\subset G_L$, $p(G_L)=H$, and $L/K$ is finite and
contains $K_1$. Its residue field is the field fixed by $H$. Field base
change gives continuous extension isomorphisms
\[
 E_{U_L}=E_U\times_{G_K}G_L,\qquad
 E_{U_L}\times_{G_L,s|_H}H=E_s\times_{G_k}H.
\]
Thus a nodal stable model after finite extension obstructs a proposed smooth
realization after restricting to this particular open subgroup. Stable
models exist after finite extension by properness of the stable pointed
moduli stack, and a node persists under every further field extension of
the stable special fibre \citep[Theorem 2.7]{BC_Knudsen}.

Conversely take a finite Galois extension $K_1/K$ giving good marked
reduction, with group $G$ and residue field $k_1$. Uniqueness of stable
models extends every generic semilinear $G$-isomorphism to the model.
Write $\lambda_\sigma$ for the reduced semilinear isomorphism of its special
pair $(C_1,D_1)$. The map $\alpha_s:G_k\to G_K\to G$ is continuous and
has finite image. Over the finite Galois residue extension fixed by
$\ker\alpha_s$, the maps $\lambda_{\alpha_s(g)}$ give an actual descent
datum. It is effective: a sufficiently high power of
$\omega_{C_1}(D_1)$ embeds the pair, its section space is a finite-dimensional
semilinear vector space, and ordinary Galois descent descends both its
homogeneous ideal and the marked divisor. Smoothness and the etale divisor
property descend faithfully flatly. Call the descended complement $V_s$.

The comparison is equivariant for
$Q=G\times_{\Gal(k_1/k)}G_k$. The natural surjection $G_K\to Q$ has kernel
$I_{K_1}$, which acts trivially by good reduction. Restriction along $s$ is
$g\mapsto(\alpha_s(g),g)$, exactly the descended action. Notice that the
second coordinate retains the entire arithmetic action on geometric
fundamental groups; this action is not asserted to factor through a finite
group. B.2 now identifies the extensions. This proves all three
equivalences for the specified $s$.
\end{proof}

\subsection{B1--B2: twists, homology and weights}

\begin{claim}{B.5}{Known}
For a split stable marked model, positive Kummer inertia acts by the
profinite multi-twist $\prod_e\tau_e^{m_e}$, and its map from $T$ is
injective if a node exists. On homology it acts by
$x\mapsto x+\sum_e m_e\langle x,\delta_e\rangle\delta_e$.
Exchanging branches inverts the oriented edge generator; it does not insert
an additional minus sign in the unoriented Dehn twist.
\end{claim}
\begin{proof}
The input is \citet[Theorem 4.8(iv)--(v), Lemma 5.4(ii)--(iii),
Corollary 5.9]{BC_HMTopics}: for a semi-graph of pro-$\Sigma$
pointed-stable-curve type its Dehn group is the product of one cyclotome per
node; the log deformation map has coordinates the lengths $m_e$; conjugation
permutes those coordinates and multiplies them by the cyclotomic character.
These statements hold for every nonempty set of primes excluding the
residue characteristic. Here that is all primes. The chart $xy=\varpi^m$
of B.3 fixes the positive Kummer normalization; pulling a loop in the
smoothing parameter back under $t_e=\varpi^{m_e}$ multiplies its exponent
by $m_e$. The topological sign agrees with the plumbing convention
$y=\varpi^m/x$. Across a separating edge it conjugates the second vertex
group by $\gamma_e^{m_e}$; across an HNN edge it sends the stable letter
$u$ to $u\gamma_e^{m_e}$. These are identities on finite coverings and
therefore pass to the profinite groups. The corresponding topological
formula is also explicitly Theorem (3.2) of
\citet{BC_AMO}.

Changing paths conjugates the construction by the geometric group and
therefore does not change its outer class. Reversing an edge replaces its
generator by the inverse and also reverses the crossing direction; the
two changes cancel in the twist. If a node is present, the coordinate map
$T\to T$, $a\mapsto m_ea$, is injective, proving inertia injectivity.
On a crossing cycle the twist adds $m_e$ copies of the vanishing circle;
signed crossing number is $\langle x,\delta_e\rangle$. The circles are
disjoint, so their pairwise intersections vanish and the contributions add.
This derives the abelian formula from the actual group automorphism.
In the nonsplit case Galois permutes nodes and multiplies an oriented edge
by $\chi$ times the sign of its branch permutation. In the homology formula
both occurrences of $\delta_e$ change sign, while in the Dehn-coordinate
formula only $\chi$ and the node permutation remain.
\end{proof}

\begin{claim}{B.6}{Proved}
The filtration and dimensions in B2/L2 hold. In the nonsplit version the
weight-zero piece is $H_1(\mathcal G,\Q_\ell)$ with its graph action, the
weight-minus-one piece is $\bigoplus_vV_\ell\Jac(C_v)$ with component
permutations and descent action, and the weight-minus-two piece is the
module of boundary cycles of the cut surface, with branch signs and Tate
twist. There is a natural exact sequence
\[
 0\longrightarrow
 (\Q_\ell[D(\bar k)]/\Q_\ell)(1)
 \longrightarrow W_{-2}\longrightarrow H^1(\mathcal G,\Q_\ell)(1)
 \longrightarrow0
\]
when $\nu>0$; for $\nu=0$ the first term is zero. The phrase ``node part''
means the image of monodromy, a complement over $\Q_\ell$, rather than an
unspecified canonical direct summand of the integral boundary lattice.
\end{claim}
\begin{proof}
Cut the smooth surface at its vanishing circles and retain the actual
punctures. Boundary homology of the resulting components, modulo opposite
branch identifications and the sum-of-boundaries relation at every vertex,
is $W_{-2}$. Filling those boundaries gives the homology of the normalized
components; killing all vertex homology gives graph homology. The cellular
chain complex of this decomposition proves the three successive quotients.
Counting relations in the connected graph gives respectively
$t+\max(\nu-1,0)$, $2\sum_vg_v=2(g-t)$, and $t$.
Every construction uses boundary or vertex maps of the semi-graph, hence
is Galois-equivariant in the log comparison B.3.

Log inertia acts trivially on these quotients: edge twists change a
crossing by a boundary class and act trivially on the normalized components.
Consequently their $G_k$ actions are independent of the residue splitting.
For an oriented cycle $a=(a_e)$ in the graph the monodromy pairing is
$q(a,b)=\sum_em_ea_eb_e$. Its matrix on a cycle basis has positive
determinant: over $\R$ it is positive definite, since all $m_e>0$.
Its determinant is a nonzero integer, hence nonzero in every $\Q_\ell$.
The Picard--Lefschetz formula of B.5 therefore identifies the graph quotient
with the $t$-dimensional image of $N=\log\tau=\tau-1$ in $W_{-2}$.
The quotient of boundary cycles by the cusp subspace is the dual graph
homology with Tate twist, giving the displayed exact sequence.

There is a genuine open-curve caution here. If ``node part'' were read as
the span of all individual vanishing-cycle classes, that strengthened
reading would be false. For $\mathbf P^1\setminus\{0,t,1,\infty\}$ the
stable graph is a tree, so $t=0$ and $N=0$, but the separating node has
nonzero homology class $a_0+a_t$ in the cusp augmentation module. It is in
the radical of the intersection form. Thus the isomorphism asserted for
$N$ is to its image, which projects isomorphically onto the dual graph
quotient, and not to the span of every boundary circle called a node.

Finally choose a finite residue extension splitting the finite graph and
marks and models of all normalized component Jacobians. Specialize those
models at a common good finite-field point. Frobenius on the three graded
pieces has complex absolute values $1$, $q^{1/2}$ and $q$, respectively,
with the convention of arithmetic Frobenius on homology. Thus their
characteristic polynomials have disjoint roots. A shared constituent would
give a common Frobenius eigenvalue on the corresponding subquotient, a
contradiction. The same argument works after any prescribed finite
extension of $k$. The statement is about constituents, not about splitting
the possibly nontrivial extensions between the pieces.
\end{proof}

\subsection{B3--B4: the exact abelian criterion}

The margin of cote 149, manuscript p.~49, scan 100, says ``je l'ai
d\'emontr\'e, dans le contexte des VA, il y a belle lurette''. The statement
that the opened SGA 7 text supplies is the abelian-variety criterion below.
The computations distinguish that criterion from good reduction of the
curve; the handwriting alone does not certify a stronger historical claim.

\begin{claim}{B.7}{Proved}
For any henselian trait of residue characteristic zero, the following are
equivalent: inertia is trivial on $H_1(U_{\bar K},\Z_\ell)$ for one prime
$\ell$; it is trivial on $\Pi^{\rm ab}$; the Jacobian of $\bar U$ has good
reduction and inertia fixes every geometric point of $D$. Equivalently the
generalized Jacobian extends as an extension of an abelian scheme by a
torus. The same statement about one prime $\ell\ne p$ and all prime-to-$p$
homology holds in residue characteristic $p$.
\end{claim}
\begin{proof}
The generalized Jacobian $J_D$ fits into $0\to S_D\to J_D\to J\to0$,
where $S_D=(\operatorname{Res}_{D/K}\G_m)/\G_m$ if $D\ne\varnothing$ and
$S_D=0$ otherwise. Its Tate module is $H_1(U_{\bar K},\Z_\ell)$; the
torus character lattice is the degree-zero lattice on $D(\bar K)$.
The resulting exact sequence is
\[
0\to(\Z_\ell[D(\bar K)]/\Z_\ell)(1)\to T_\ell J_D
\to T_\ell J\to0.
\]
For $\ell$ invertible in $R$, inertia acts trivially on $\Z_\ell(1)$.
Its permutation action on the displayed quotient is faithful when
$\nu\ge2$; for $\nu=1$ the one point is automatically fixed, and for
$\nu=0$ there is no condition. Trivial action on $T_\ell J_D$ therefore
implies unramified marks and trivial action on $T_\ell J$.

SGA 7 I, IX, Corollary 5.10 says that an abelian variety has good reduction
if and only if its $\ell$-divisible group extends over the trait, assuming
$\operatorname{char}K=0$ when $\ell=p$. For $\ell\ne p$ an extension is
equivalent to an unramified Tate module, since finite etale covers of a
henselian trait are precisely unramified Galois modules. This gives the
claimed good reduction of $J$ \citep[IX, Corollary 5.10 and
Proposition 5.13(a)]{BC_SGA7}.

Conversely, unramified marks extend $S_D$ to a torus over $R$. After a finite
unramified splitting extension, an extension of $J$ by $S_D$ is specified
by the homomorphism from the character lattice of $S_D$ to $J^\vee(K)$.
Every value extends uniquely to an $R$-section of the proper abelian scheme
$J^\vee$. The Poincare biextension then constructs the extended group
extension; uniqueness and the original Galois action descend it over $R$.
Its $\ell^n$-torsion is finite etale for $\ell$ invertible in $R$, proving
unramifiedness of every such Tate module. In residue characteristic zero
this includes every prime, so the product is $\Pi^{\rm ab}$. This also
shows why collisions of distinct marked sections in a chosen curve model
do not obstruct this abelian criterion: their divisor classes still extend
to the proper dual abelian scheme.
\end{proof}

\begin{claim}{B.8}{Proved}
For residue field finitely generated over $\Q$, the complete implication
table for the five conditions of scans 100--103 is as follows. A row is the
antecedent; a column is the consequent.
\[
\begin{array}{c|ccccc}
 &a&b&c&d&e\\\hline
a&-&\mathrm{yes}&\mathrm{yes}&\mathrm{yes}&\mathrm{yes}\\
b&\mathrm{yes}&-&\mathrm{yes}&\mathrm{yes}&\mathrm{yes}\\
c&\mathrm{no}&\mathrm{no}&-&\mathrm{yes}&\mathrm{no}\\
d&\mathrm{no}&\mathrm{no}&\mathrm{yes}&-&\mathrm{no}\\
e&\mathrm{no}&\mathrm{no}&\mathrm{no}&\mathrm{no}&-
\end{array}
\]
\end{claim}
\begin{proof}
Tamagawa's Theorem (5.3), for a DVR and a hyperbolic smooth marked generic
curve, equates good marked reduction with trivial prime-to-residue-characteristic
outer inertia, and with trivial outer pro-$\ell$ inertia for one $\ell$
\citep[Theorem (5.3)]{BC_Tamagawa}. In residue characteristic zero this is
$a\Leftrightarrow b$. Inner automorphisms act trivially on abelianization,
so the characteristic quotient maps give $b\Rightarrow c\Rightarrow d$;
B.7 gives $d\Rightarrow c$. Good reduction gives the special-fibre action,
hence $a\Rightarrow e$.

For all negative implications from $c$ or $d$, use
$\mathbf P^1\setminus\{0,t,1,\infty\}$ over $\Q((t))$. Blowing up the
meeting sections $0,t$ gives two rational components, each with three special
points, joined at a node of thickness one. The nontrivial Dehn twist of B.5
is a separating twist; its vanishing class is a sum of cusp classes and is
in the radical of intersection pairing, so homological inertia is trivial.
B.4 excludes $e$. The model is stable nodal, so $a,b$ fail.

For negative implications from $e$, use
$y^2=x^3-t^2x$ minus the origin over $\Q((t))$. With $r^2=t$, put
$x=r^2X$, $y=r^3Y$; the equation becomes $Y^2=X^3-X$.
The nontrivial inertia element acts as $[-1]$ on the elliptic curve and
as $-1$ on its Tate module. Thus $b,c,d$ fail and $a$ fails by Tamagawa.
Good reduction after adjoining $r$ and B.4 give $e$ for every splitting.
The coordinate substitutions, discriminants and homology matrices are
checked in \texttt{code/BC\_checks.py}.
\end{proof}

\begin{claim}{B.8u}{Proved}
With the full geometric group in $b,c,e$ and $\ell\ne p$ in $d$, the
table of universal implications over henselian traits of positive residue
characteristic is
\[
\begin{array}{c|ccccc}
 &a&b&c&d&e\\\hline
a&-&\mathrm{no}&\mathrm{no}&\mathrm{yes}&\mathrm{no}\\
b&\mathrm{yes}&-&\mathrm{yes}&\mathrm{yes}&\mathrm{yes}\\
c&\mathrm{no}&\mathrm{no}&-&\mathrm{yes}&\mathrm{no}\\
d&\mathrm{no}&\mathrm{no}&\mathrm{no}&-&\mathrm{no}\\
e&\mathrm{no}&\mathrm{no}&\mathrm{no}&\mathrm{no}&-
\end{array}
\]
The same table gives the universal implications when restricted to equal
characteristic $p>0$. This is a statement allowing $p$ to vary; the
counterexamples below are given at $p=3$ and $p=7$, not asserted to
occur by the same formulas at every prime. The separate
mixed-characteristic table, where $b,c,e$ have impossible antecedents,
is C.3. The prime-to-$p$ table is C.5.
\end{claim}
\begin{proof}
The implications $b\Rightarrow a$ and $a\Rightarrow d$ follow from
Tamagawa's criterion and smooth prime-to-$p$ specialization, respectively;
$b\Rightarrow c\Rightarrow d$ is functorial abelianization and projection
to the $\ell$-component. B.20 proves $b\Rightarrow e$ in equal
characteristic, by a full specialization argument for proper curves and
by the impossibility of $b$ for affine curves. In mixed characteristic
the implication is vacuous by C.3.
For all negative cells B.21 gives $a,d$ true and $b,c,e$ false;
B.22 gives $c,d$ true and $a,b,e$ false; and B.23 gives $e$ true and
$a,b,c,d$ false. These exhaust the twenty ordered pairs. Constant proper
hyperbolic curves show that $b$ is possible in equal characteristic,
so its positive row is not merely vacuous there.
\end{proof}

\subsection{B5--B6: what varies with the splitting}

\begin{claim}{B.9}{Refuted}
The splittings themselves do not form a torsor under $H^1(k,T)$ as L4 says.
After a reference splitting $s_0$ is fixed, they form a torsor under
$Z^1_{\rm cts}(G_k,T)$; their $T$-conjugacy classes form a torsor under
$H^1(k,T)$. If the outer inertia image has finite order $n$, good marked
reduction holds over $K(\varpi^{1/n})$, and the realizing curve $V_s$ is the
twist of $V_{s_0}$ by the image of this cohomology class under
$H^1(k,T)\to H^1(k,\mu_n)\to H^1(k,\Aut(C',D'))$.
\end{claim}
\begin{proof}
Write $G_K=T\rtimes G_k$ using $s_0$, and write
$s(g)=(a(g),g)$. Multiplication gives
$a(gh)=a(g)+\chi(g)a(h)$, exactly the continuous cocycle condition.
Conjugation by $u\in T$ replaces $a(g)$ by $a(g)+u-\chi(g)u$.
These are distinct splittings whenever the coboundary is nonzero; for
example complex conjugation over $\Q$ gives $u-\chi(g)u=2u\ne0$ for
$u\ne0$. Thus the distinction is substantive. Quotienting by these
coboundaries gives precisely the asserted $H^1$ torsor.

A finite quotient of $T=\widehat\Z(1)$ of order $n$ has kernel $nT$.
The extension obtained by adjoining an $n$th root of a uniformizer is totally
ramified of degree $n$ and has inertia $nT$ in the chosen separable closure.
Tamagawa's criterion therefore gives good marked reduction over it.
Conversely good reduction after finite extension makes an open subgroup of
$T$ act trivially, so the inertia image is finite. This proves the
finiteness criterion without replacing ``good'' by ``potentially good''.

Take $s_0$ to fix compatible roots of $\varpi$. The special pair over
$K(\varpi^{1/n})$ has a descent action of $T/nT=\mu_n$; conjugation by
$s_0(G_k)$ is the cyclotomic action, making
$\mu_n\to\Aut(C',D')$ a $G_k$-equivariant homomorphism. If the root extension
is not Galois, these statements are made after adjoining $\mu_n$ and
descending the resulting finite etale group scheme; no constant cyclic
Galois group over $k$ is silently assumed. The finite descent datum in B.4
changes by exactly the image of $a(g)$ in this automorphism group. The
standard cocycle construction of a twist gives the displayed formula.
Changing $a$ by a coboundary gives an isomorphic descended curve, as required.
\end{proof}

\begin{claim}{B.10}{Proved}
No isomorphism-invariant predicate on the residue action alone can be
equivalent to good marked reduction. A predicate equivalent to $e$ is that
there are precisely $\nu$ geometric conjugacy classes of maximal
essentially cyclotomically normalized procyclic subgroups. The number $\nu$
can itself be recovered from the action, so this is an intrinsic predicate
on the stated class of actions. Adding triviality
of the original inertia action in $\Out(\Pi)$ makes it equivalent to $a$.
\end{claim}
\begin{proof}
For $y^2=x^3-t^2x$ minus the origin and the reference splitting in B.9,
the residue action is that of the constant once-punctured elliptic curve
$Y^2=X^3-X$ over $\Q$. The first family has bad marked reduction and the
constant family has good reduction. Thus the residue action cannot
distinguish them. For genuine point actions choose $(0,0)$ on both affine
elliptic curves: the coordinate isomorphism and its $[-1]$ descent datum
fix this point, so their residue point actions agree as well.
The count criterion follows from B.4: after the explicit restriction
comparison there, a stable nodal model contributes at least $\nu+n$ such
classes, while a smooth model has exactly $\nu$. Essential normalization
is unchanged by passing to an open subgroup: intersect a normalizing group
with its inverse image for one direction, and use the same normalizing
group for the other. This proves the reduction to a stable model and covers
the proper case $\nu=0$.

To make the predicate intrinsic, put $M=H_1(\Pi,\Q_\ell)$. If
$H^2(\Pi,\Z/\ell)\ne0$, set $\nu(\rho)=0$. Otherwise let
\[
 a=\max_{k'/k\ {
m finite}}\dim\Hom_{G_{k'}}(M,\Q_\ell),\qquad
 b=\max_{k'/k\ {
m finite}}\dim\Hom_{G_{k'}}(\Q_\ell(1),M).
\]
These maxima exist since the dimensions are bounded integers and both
maxima can be attained after one common finite extension. Use B.4 to reach
a stable model compatibly with the given splitting and then split its
graph and marks by a further residue extension. B.1 and B.6 give
$a=t$ and $b=t+\nu-1$: every map to the trivial module kills $W_{-1}$,
and every map from $\Q_\ell(1)$ lands in $W_{-2}$. The same vanishing
holds after every further finite extension. Hence
$\nu(\rho)=b-a+1$ in the affine case. The predicate is that the above
intrinsic number of cyclotomic maximal classes equals $\nu(\rho)$.
Finally the added inertia condition is exactly
Tamagawa's criterion, and good reduction already implies $e$.
\end{proof}

\subsection{B7: the residue-prime warning}

\begin{claim}{B.11}{Known}
Over a henselian mixed-characteristic trait an abelian variety has good
reduction if and only if its $p$-divisible group extends over the trait.
For a smooth proper genus-at-least-two curve with a semistable model over
a complete characteristic-zero DVR with perfect residue field of
characteristic $p$, and a base point extending to a section of that model,
good reduction is equivalent to crystallinity of its unipotent $p$-adic
etale fundamental group, equivalently of every $E^{\rm et}_{n,b}$.
These are distinct criteria; the first does not detect good reduction of
the curve.
\end{claim}
\begin{proof}
The first assertion is exactly SGA 7 I, IX, Corollary 5.10 and
Proposition 5.13(a), including their explicit $\operatorname{char}K=0$
hypothesis for $\ell=p$ \citep{BC_SGA7}. The second is
\citet[Theorems 1.6 and 1.9]{BC_AIK}; their introduction fixes completeness,
perfect residue field, properness, genus at least two, semistability and
an integral base point before stating these theorems. We retain all of those
hypotheses. The Galois exact sequence is now
$1\to I_K\to G_K\to G_k\to1$ with
$1\to P_K\to I_K\to\prod_{\ell\ne p}\Z_\ell(1)\to1$.
Replacing $I_K$ by the full cyclotomic $T$ is false.

The example $y^2=(x^3-5^6)(x^3-1)$ over $\Q_5$ has stable special fibre two
smooth elliptic curves meeting in a separating node of thickness one.
Indeed separate the branch clusters by $xz=5^2$ and set
$A=(1-z^3)(x^3-1)$; at the join normalization has
$V^2=xA$, $W^2=z/A$, $VW=5$. The unit $A(0,0)=-1$ and the implicit function
criterion give the regular nodal chart. The two components have equations
$Y^2=1-X^3$ and $V^2=x(x^3-1)$, smooth in characteristic five. Thus the
Separating-node Picard--Lefschetz gives trivial prime-to-five homological
inertia, so the abelian criterion B.7 proves that the Jacobian has good
reduction. Equivalently the relative degree-zero Picard
scheme of a compact-type stable curve is an abelian scheme, with special
fibre the product of the component Jacobians
\citep[Introduction, compact-type models]{BC_Dupuy}. Its $5$-divisible
group therefore extends. The nodal fibre is stable,
so the curve has bad reduction. At either rational point at infinity the
proper model supplies an integral base point. AIK therefore implies that
at least one $E^{\rm et}_{n,b}$ is not crystalline. The polynomial and local
chart identities are checked in \texttt{BC\_checks.py}.
\end{proof}

\begin{claim}{B.11a}{Proved}
Let $K$ be a complete discretely valued field of characteristic zero with
perfect residue field of characteristic $p$, and let $C/K$ be a smooth
proper geometrically connected curve of genus at least two with $b\in C(K)$.
Then $C$ has good reduction if and only if every representation
$E^{\rm et}_{n,b}$ of its pointed unipotent $p$-adic fundamental group is
crystalline. No semistability assumption is necessary in this proper
statement.
\end{claim}
\begin{proof}
The universal unipotent objects and their fibres are defined from
$\pi_1(C_{\bar K},b)$ independently of a model. The first nonconstant
graded piece, at $n=2$, is dual to $H^1_{\rm et}(C_{\bar K},\Q_p)$;
this is the identification appearing in \citet[Theorem 1.7(ii)]{BC_AIK}.
Crystalline representations are closed under subquotients and duals, so
crystallinity of every $E^{\rm et}_{n,b}$ implies crystallinity of
$V_p\Jac(C)$. Their Theorem 1.2 then gives good reduction of the Jacobian.

Deligne--Mumford's Theorem 2.4 says that a proper genus-at-least-two curve
has stable reduction if and only if its Jacobian has semiabelian reduction
\citep[Theorem 2.4, printed p.~89]{BC_DM}. Its Section 2 assumes the residue
field algebraically closed. To apply it here, pass to the strict
henselization of $R$; the perfect-residue-field hypothesis makes its residue
field algebraically closed. Good reduction of the Jacobian is preserved
by this base change. The unique stable model supplied there descends:
it and its isomorphisms are of finite presentation, hence defined over a
finite unramified extension, and uniqueness of stable models makes the
generic Galois descent datum extend and satisfy its cocycle identity.
Effective descent for the canonically polarized model gives a stable model
over $R$. This is precisely the descent mechanism in the proof of the
same theorem on printed p.~92. Its resolution has regular semistable total
space, by Proposition 2.3 of that paper. Properness extends $b$ to an
$R$-point of this model. All hypotheses of AIK's Theorem 1.9 now hold, so
$C$ has good reduction. Conversely, good reduction supplies such a proper
smooth model and integral base point, and AIK's theorem makes every
$E^{\rm et}_{n,b}$ crystalline.
\end{proof}

\begin{claim}{B.11b}{Proved}
For a pair $(E,\{O\})$ of type $(1,1)$ over the complete mixed-characteristic
field with perfect residue field of B.11a, good marked reduction is
equivalent to crystallinity of
$H_1((E\setminus\{O\})_{\bar K},\Q_p)$, with no semistability assumption.
\end{claim}
\begin{proof}
With one puncture the boundary augmentation module is zero, so this
representation is $V_pE$. AIK's Theorem 1.2 equates its crystallinity with
good reduction of $E$. The identity section extends over the abelian scheme
and is a relative etale divisor, hence gives good marked reduction.
Conversely a good marked model gives good reduction of $E$. In the local
field case the same relative criterion is explicitly recorded in
\citet[Theorem 3.1 and Introduction]{BC_Porowski}; the additional $p\ge5$
and torsion/discrete-section assumptions of his absolute reconstruction
theorems are not needed for this relative assertion.
\end{proof}

\begin{claim}{B.11c}{Known}
For a stable marked log curve over a standard log point over a
characteristic-zero field, its unipotent log de Rham fundamental group's
outer monodromy is trivial if and only if the underlying stable curve is
smooth. A section specializing to the smooth unmarked locus gives trivial
genuine monodromy in the smooth case. This is a marked log de Rham theorem,
not by itself a mixed-characteristic etale crystalline comparison.
\end{claim}
\begin{proof}
These are exactly \citet[Theorem 6.5 and Proposition 6.8]{BC_CDPS};
Definition 6.4 includes marked stable curves, and Definition 6.6 defines
the good-section condition as lying in the relative trivial-log locus.
Their Theorem 6.13 proves the nontrivial outer monodromy in the minimal
semistable case. The mixed-characteristic application in their Theorem 6.29
explicitly assumes a proper curve and invokes AIK's comparison. It does
not state the additional marked etale comparison required for all of B7.
\end{proof}

\begin{claim}{B.11d}{Proved}
Let $K$ be complete discretely valued of characteristic zero, with perfect
residue field of characteristic $p>0$, and let $U=\bar U\setminus D$ be
any hyperbolic curve. Good marked reduction is equivalent to the following
criterion in $p$-adic data of finite covers of $\Pi_U$:

There exist a finite \emph{unramified} extension $L/K$, a geometrically
connected finite Galois cover $V\to U_L$ of degree prime to $p$, and
$v\in V(L)$ such that its smooth proper compactification $C$ has genus
at least two, the finite map $C\to\bar U_L$ is ramified at every
geometric cusp, and the proper unipotent $p$-adic fundamental group of
$C$ based at $v$ is crystalline. Equivalently, every universal unipotent
representation $E^{\mathrm{et}}_{n,v}$ of $G_L$ for $C$ is crystalline.

In the forward direction the cover can have degree a power of any fixed
prime $\ell\ge5$, $\ell\ne p$. For $D=\varnothing$ the identity cover
is allowed. No semistability assumption on $U$ or $C$ is needed.
\end{claim}
\begin{proof}
We first identify the data attached to the original group. Such a cover
is given by a continuous quotient $q:E_{U_L}\twoheadrightarrow Q$
whose restriction to $\Pi_U$ is surjective. Put $H=\ker(q|_{\Pi_U})$.
The proper geometric group of its compactification is
\[
 \Pi_C=H\big/\overline{\langle H\cap\gamma I_c\gamma^{-1}:
                 c,\ \gamma\in\Pi_U\rangle}^{\mathrm{normal}},
                                                               \tag{B.11d.1}
\]
where $I_c$ ranges over the geometric cusp groups. Indeed a finite
etale cover of $V$ extends etale to $C$ exactly when its inertia at
each missing point vanishes, by normalization and the local parameter
criterion. Quotienting the arithmetic subgroup $\ker q$ by the same
normal subgroup gives $E_C$. The point $v$ defines its actual Galois
action; its $\Q_p$-unipotent completion is the data in the criterion.
We use these proper quotients of finite-index subgroups of $\Pi_U$,
with inherited arithmetic action and peripheral data, rather than
claiming direct crystallinity of $\Pi_U$ at every base point.

Suppose the pair has good reduction. A finite unramified extension
splits its etale marked divisor and supplies an integral point $b$ of
the smooth open model: a nonempty smooth special fibre has a point
over a finite separable residue extension, and henselian lifting
supplies $b$. We choose a finite $\ell$-group quotient in which every
cusp survives. If $\nu\ge2$, abelianization modulo $\ell$ suffices,
since the boundary classes have their sole relation that their sum is
zero and each class is nonzero. If $\nu=1$, then $g\ge1$, and the
geometric pro-$\ell$ group is free on $a_1,b_1,\ldots,a_g,b_g$,
with cusp $(\prod_i[a_i,b_i])^{-1}$. Map $a_1,b_1$ to the two
elementary generators of the Heisenberg group over $\F_\ell$, all
other generators to one; their commutator is a nonzero central element.
These presentations are SGA 1, XIII, Corollary 2.12 \citep{BC_SGA1}.

Intersecting all automorphic images of the kernel makes it
characteristic. Only finitely many occur: open subgroups of a fixed
index in a finitely generated profinite group are determined by
homomorphisms to a fixed finite symmetric group and a point stabilizer.
The new quotient $Q$ embeds in a finite product of the old finite
$\ell$-group, so remains an $\ell$-group and no cusp becomes trivial.
Every ramification index $e_c$ is at least $\ell$. The genus of the
compactification satisfies
\[
 2g(C)-2=|Q|\left(2g-2+\sum_{c\in D}(1-e_c^{-1})\right).
                                                               \tag{B.11d.2}
\]
For $g=0$, $\nu\ge3$ and $\ell\ge5$ bound the parentheses below by
$-2+3(1-1/5)>0$; for $g=1$, the nonempty ramification set makes them
positive; for $g\ge2$ positivity is automatic. Thus $g(C)\ge2$.
When $D$ is empty we instead use the identity cover, with $g\ge2$.

The integral section $b$ splits the arithmetic geometrically pro-$\ell$
extension. Smooth-pair comparison identifies its \emph{actual} action
with the action of its special-fibre point, so inertia is trivial.
The exact inputs are SGA 1, XIII, Proposition 4.3, Examples 4.4 and
4.5.3, with $\ell$ invertible on the trait \citep{BC_SGA1}.
The finite action on the characteristic quotient $Q$ therefore factors
through $G_k$. A further finite unramified extension kills it. Projection
to $Q$ is then a homomorphism of the semidirect product, trivial on
the section of $b$, and defines the desired constant-deck Galois cover
with an $L$-point $v$ above $b$.

The same quotient of the smooth open model's arithmetic group extends
the cover etale over that model. Normalize the proper model in its
function field. Ramification occurs only along the horizontal marks,
with indices powers of $\ell$. SGA 1, XIII, Propositions 5.1--5.2
give, after an etale change, the regular Kummer chart $z=w^{e_c}$.
Thus the normalization is smooth over the trait at the missing points
as well. The proper curve $C$ has good reduction, and $v$ is integral;
AIK supplies the crystalline representations
\citep[Theorems 1.6 and 1.9]{BC_AIK}.

Conversely suppose the criterion holds. B.11a gives good reduction
of the proper curve $C$, without semistability in advance. Its deck
automorphisms extend to its smooth proper model $\mathcal C/R_L$,
and their special action remains faithful by the finite unramified
stable-curve Isom scheme \citep[Theorem 1.11]{BC_DM}.
The projective quotient $\mathcal X=\mathcal C/Q$ exists by the
invariant graded ring of the canonical polarization. We prove locally
that it is smooth and that its reduced branch divisor is finite etale.

After strict henselization and completion, take a geometric special
point with stabilizer $H\subset Q$. Its tangent character is faithful.
Indeed a nontrivial finite-order automorphism of $\bar k[[z]]$ whose
order is prime to $p$ cannot have derivative one: its first nonzero
higher coefficient is multiplied by that order under iteration, and
cannot vanish. An automorphism identical on the completed local ring
is identical on the curve, contrary to faithfulness of the special
action. Thus $H$ embeds in $\bar k^*$ and is cyclic, of order $e$.

If $e>1$, a generator $h$ has special derivative $\zeta\ne1$.
The formal equation $h(z)-z=0$ has invertible derivative, so Hensel's
lemma gives a unique fixed section lifting the special point. Shift
$z$ to vanish on that section, lift $\zeta$ to the $e$th roots of
unity in the trait, and put
\[
              w=e^{-1}\sum_{i=0}^{e-1}\zeta^{-i}h^i(z).
\]
Its special linear coefficient is nonzero and $h(w)=\zeta w$.
The completed invariant ring is consequently
$R_L[[w]]^H=R_L[[w^e]]$. For $H=1$ the ring is simply $R_L[[w]]$.
To obtain the local ring of the global quotient, the completed
semilocal ring of a $Q$-orbit is a product of these local rings;
taking $Q$-invariants reduces it to $H$-invariants in one factor.
This proves smoothness of $\mathcal X$. At a ramified point the
branch locus is exactly the section $w^e=0$. No branch point appears
only in the special fibre, because every special fixed point lifted
uniquely, and no two branch points merge in this chart. Properness
then makes the full reduced branch divisor finite etale. Smoothness and this
divisor property descend through the strictly local base changes.

The generic quotient is $\bar U_L$, and its reduced branch divisor is
exactly $D_L$ by hypothesis. Thus $\mathcal X$ with that divisor is a
good marked model. Finally $L/K$ is unramified, so $I_L=I_K$ under
the usual common separable closure identification. Good marked
reduction over $L$ makes the pro-$\ell$ outer inertia trivial already
over $K$; Tamagawa's criterion gives good marked reduction over $R$
\citep[Theorem 5.3]{BC_Tamagawa}.
\end{proof}

This treats every marked type, including $(0,3)$, $(0,4)$ and $(1,1)$,
under the complete/perfect-residue hypotheses, through proper AIK on
a finite cover. A direct marked comparison on the original open
curve's unipotent group is not presumed. B.11f--B.12 give a separate
finite-cover $p$-divisible-group criterion over arbitrary residue fields.

\begin{claim}{B.11e}{Proved}
Let $R$ be henselian, of mixed characteristic, with perfect residue
field, and let $\widehat K$ be the fraction field of its completion.
Good marked reduction over $R$ is equivalent to the criterion B.11d
applied to $U_{\widehat K}$. Its Galois representations are attached
to the original arithmetic group via
\[
 G_{\widehat K}\simeq G_K,\qquad
 E_{U_{\widehat K}}\simeq E_U.
\]
In this formulation the auxiliary points in B.11d are points over
finite unramified extensions of $\widehat K$; no descent of those
points to $K$ is asserted.
\end{claim}
\begin{proof}
The natural continuous outer map $G_{\widehat K}\to G_K$ is an
isomorphism by the Krasner comparison, in exactly the henselian-DVR
scope of \citet[Lemma 3.1]{BC_MinTsuj2022}. Its map on residue Galois
groups is the identity, so it identifies inertia. Characteristic-zero
geometric base change identifies the geometric curve groups, including
the cusp subgroups, and the two arithmetic exact sequences therefore
identify $E_{U_{\widehat K}}$ with $E_U$.
Good marked reduction over $R$ implies it over $\widehat R$, hence
the finite-cover crystalline criterion by B.11d. Conversely that
criterion implies good marked reduction over $\widehat R$.
Thus the pro-$\ell$ outer inertia action is trivial there, for
$\ell\ne p$, and the identified inertia action is trivial over $K$.
Tamagawa's criterion over the original DVR gives good marked reduction
over $R$. This last argument avoids any unproved algebraization from
a formal model and retains the perfect residue field needed by AIK.
\end{proof}


\begin{claim}{B.11f}{Proved}
Let $R$ be a henselian DVR with fraction field $K$ of characteristic
zero and residue field of characteristic $p>0$, with no perfection
assumption. Let $C/K$ be a smooth proper geometrically connected curve
of genus $h\ge2$, and fix one prime $\ell\ne p$. Define
$\mathcal P_\ell(C/R)$ by the following finite collection of conditions:

\begin{enumerate}
\item $\Jac(C)[p^\infty]$ extends to a $p$-divisible group over $R$.
\item For each connected geometric cyclic etale cover
$D\to C_{\bar K}$ of degree $\ell$, there is a finite extension
$L/K$ and a descent $D_L\to C_L$ such that
$\Jac(D_L)[p^\infty]$ extends to a $p$-divisible group over $R_L$.
\end{enumerate}
Here $R_L$ is the valuation ring of the unique extension of the
valuation. The extensions in the second condition may be ramified.
There are exactly $(\ell^{2h}-1)/(\ell-1)$ isomorphism classes of
these unlabelled geometric covers, so one common finite extension
can be used for all of them. Then
\[
          C\text{ has good reduction over }R
          \quad\Longleftrightarrow\quad\mathcal P_\ell(C/R).
\]
The conditions are conditions on integral $p$-adic Galois data of
$\Pi_C$ and its cyclic index-$\ell$ subgroups: their abelianized
pro-$p$ groups are the corresponding Jacobian Tate modules.
\end{claim}
\begin{proof}
In characteristic zero a $p$-divisible group over a field is an etale
system of finite group schemes, determined by its integral Tate module.
For the Jacobian this module is
$H_1(\Pi_C,\Z_p)$. Condition~(1) therefore asks for an integral
Barsotti--Tate extension of this representation. Condition~(2) asks
for potentially such extensions for a finite collection of open
geometric subgroups, with the arithmetic actions obtained after
descent of the covers. SGA~7~I, IX, Corollary~5.10 and
Proposition~5.13(a)--(b) equate these conditions respectively with good
and potentially good reduction of the Jacobians; at the residue prime
their hypothesis is $\operatorname{char}K=0$, without a perfect-residue
assumption \citep{BC_SGA7}. Different finite-field descents of a fixed
geometric cover become isomorphic over a further finite extension.
Thus the potential condition is independent of the descent chosen.

Connected cyclic degree-$\ell$ covers correspond to nonzero continuous
characters $\Pi_C\to\F_\ell$, up to multiplication by
$\F_\ell^*$. Since $H^1(C_{\bar K},\F_\ell)$ has dimension $2h$,
this gives the asserted finite count. Finite presentation descends
each cover and its deck action to a finite extension of $K$; taking
the compositum of the finitely many extensions proves the common-field
assertion. Inner automorphisms act trivially on their abelianizations,
so the resulting Tate-module actions do not depend on a point section.

We may prove the geometric assertion after completing $R$.
The natural $G_{\widehat K}\to G_K$ is an isomorphism by
\citet[Lemma 3.1]{BC_MinTsuj2022}, and it identifies inertia and
finite extensions. Characteristic-zero geometric comparison then
identifies the arithmetic fundamental groups, their finite etale
covers and the representations just described. Good reduction of a
Jacobian, and potential good reduction after a finite extension, are
preserved in both directions by this comparison: use the same inertia
action on a Tate module at an auxiliary prime different from $p$ and
the SGA~7 criterion. Good reduction of $C$ is likewise preserved in
both directions by Tamagawa's pro-$r$ outer-inertia criterion for
$r\ne p$ \citep[Theorem 5.3]{BC_Tamagawa}. Thus both sides of the
claimed equivalence are unchanged. This reduction ensures that the
proper lifting theorem SGA~1, X, Theorem~2.1 is used with its literal
\emph{complete local base} hypothesis.

If $C$ is good, its Jacobian is an abelian scheme, proving~(1).
For~(2), prime-to-$p$ specialization for a smooth proper curve
identifies each geometric degree-$\ell$ cover with a cover of the
geometric special fibre. That cover descends over a finite separable
residue extension and lifts over the corresponding finite unramified
extension of the complete trait. The exact inputs are SGA~1, X,
Corollary~3.9 and Theorem~2.1 \citep{BC_SGA1}.
Its lifted proper etale cover is smooth, so its Jacobian is good.
This proves the required potential condition for every geometric cover.

Conversely assume~(1) and~(2). By~(1) the Jacobian is good over $R$.
Take a finite extension over which $C$ has a stable model; stable
reduction applies over an arbitrary residue field, by properness of
the stable-curve stack \citep[Theorem 2.7]{BC_Knudsen}.
Pass to a further finite unramified extension when splitting geometric
components and nodes is necessary. Work on the geometric special fibre
$C_0$ and its connected dual graph $\Gamma$.
The toric rank of the Jacobian's Neron identity component is
$b_1(\Gamma)$. Precisely, SGA~7~I, IX, (12.3.6)--(12.3.7) and
(12.4.1) identify its character lattice with the kernel of the
oriented edge-to-vertex incidence map \citep{BC_SGA7}.
If the total stable model has a singularity $xy=\varpi^m$, its
regular resolution only subdivides that edge, leaving this kernel
unchanged; the same Picard--Neron comparison applies. Since the
Jacobian remains good under this base change, $b_1(\Gamma)=0$.
Thus $\Gamma$ is a tree.

Suppose there is a node $e$. Removing $e$ partitions $\Gamma$ into
two nonempty subtrees. Both sides contain a component of positive
genus. Indeed, if one side had $r$ rational components, its total
number of incident branches, including the one attachment at $e$,
would be $2(r-1)+1=2r-1$. Stability of every unmarked rational
component would require at least $3r$, a contradiction.
Choose one positive-genus component on each side and a nonzero
$\F_\ell$ character of its geometric fundamental group.

There is an etale $\F_\ell$-torsor $D_0\to C_0$ with these
characters on the two chosen components and zero characters on the
others. To verify this gluing assertion, normalization gives the exact
sequence of etale constant sheaves
\[
 0\longrightarrow\F_{\ell,C_0}\longrightarrow
 \nu_*\F_{\ell,\widetilde C_0}\longrightarrow
 \bigoplus_e\F_{\ell,e}\longrightarrow0.
\]
At a node the last map is the difference of the two branch values.
The skyscraper terms have no positive cohomology over the algebraically
closed residue field. The associated sequence therefore surjects
$H^1(C_0,\F_\ell)$ onto
$\bigoplus_v H^1(C_v,\F_\ell)$, with kernel
$H^1(\Gamma,\F_\ell)=0$. This constructs the desired class.
Its restriction to each side of $e$ is nonzero, so its cover of that
connected side is connected: a torsor under a group of prime order
is either connected or trivial. Above $e$ there are exactly $\ell$
nodes. Contracting the two connected lifted side graphs leaves two
vertices joined by $\ell$ edges. Consequently
\[
                     b_1(\Gamma_{D_0})\ge\ell-1>0.
\tag{B.pdiv.1}
\]

All the data of this finite etale torsor descend to a finite
\emph{separable} extension of the residue field: SGA~1, IX,
Theorem~4.10 makes finite etale covers invariant under a purely
inseparable extension, and finite presentation then descends them
from the separable closure \citep{BC_SGA1}. After the corresponding
unramified extension of the already chosen complete stable trait,
SGA~1, X, Theorem~2.1 lifts the cover and its $\F_\ell$-torsor
structure uniquely over the proper stable model. Its generic fibre
is geometrically connected: after any finite extension a disconnected
generic cover would disconnect the normal finite etale total cover,
whereas its connected geometric special fibre forbids that.
The lifted model is again stable, since an etale cover preserves
the nodal local equations and ampleness of the relative dualizing
bundle. It has the graph in (B.pdiv.1).

The geometric generic cover just obtained is one of those in~(2).
But its Jacobian has positive toric rank, and this remains true after
every finite field extension. Indeed the base-changed stable model
has the same geometric graph; resolving thicker nodes only subdivides
edges. The character-lattice formula therefore retains positive rank.
It cannot become an abelian scheme. A descent of this same geometric
cover satisfying~(2) would become isomorphic to it over a common finite
extension, contradicting this persistent positive toric rank. Thus
$C_0$ has no nodes, and $C$ has potentially good reduction.

It remains to descend good reduction, rather than merely assert
potential goodness. Choose a finite Galois stable extension and its
now smooth model. The finite stable-curve Isom scheme extends the
semilinear descent automorphisms, and specialization identifies their
actions on the prime-to-$p$ Tate module of the Jacobian
\citep[Theorem 1.11]{BC_DM}. By~(1), the original inertia acts
trivially on that Tate module. Each resulting special-curve
automorphism is therefore the identity. One can use
\citet[Theorem 1.13, Lemma 1.14]{BC_DM} directly: for a nonidentity
automorphism of a smooth curve of genus $h\ge2$, the graph has
nonnegative finite intersection with the diagonal, whereas trivial
action on $H^1$ would give intersection number $2-2h<0$ by their
Lefschetz--Weil calculation. This argument also covers wild finite
automorphisms.
Smooth specialization is functorial for these model automorphisms.
It follows that the entire original inertia acts trivially in the
outer group of the geometric pro-$r$ fundamental group, $r\ne p$.
Tamagawa's criterion now gives good reduction over the original trait.
This avoids an unsupported descent through a wild quotient.
\end{proof}

For the requested example $y^2=(x^3-5^6)(x^3-1)$ over $\Q_5$,
B.11 gives two elliptic special components and a separating node.
The Jacobian's $5$-divisible group extends, but a nontrivial cyclic
etale double cover on each elliptic component gives, by the preceding
construction, a geometric double cover of the genus-two generic curve
whose stable graph has two vertices and two edges. Its genus is three
and its Jacobian has toric rank one after every finite extension.
Thus condition~(2) of $\mathcal P_2$ fails, explicitly distinguishing
the curve from its good Jacobian. The graph and genus calculation is
included in \texttt{BC\_pdiv\_cover\_checks.py}.

\begin{claim}{B.12}{Proved}
Let $R$ be any henselian mixed-characteristic DVR, with residue
characteristic $p>0$, and let $U=\bar U\setminus D$ be hyperbolic.
Fix a prime $\ell\ge5$, $\ell\ne p$. Good marked reduction is
equivalent to the following criterion in $p$-adic data attached to
finite covers of $\Pi_U$:

There exist a finite \emph{unramified} extension $L/K$ and a
geometrically connected finite Galois cover $V\to U_L$ of degree
prime to $p$ whose smooth proper compactification $C$ has genus at
least two, whose extension $C\to\bar U_L$ is ramified at every
geometric cusp, and for which $\mathcal P_\ell(C/R_L)$ holds.

When $D$ is empty the identity cover is permitted. In the forward
direction the cover can have degree a power of $\ell$. No
perfection, completeness, semistability or chosen point hypothesis
is required. For perfect residue fields the crystalline finite-cover
criterion of B.11d--B.11e remains an alternative. A crystalline
comparison for the original open curve is not asserted.
\end{claim}
\begin{proof}
The proper group $\Pi_C$ is the quotient of the geometric subgroup
of the cover by the closed normal subgroup generated by its cusp
inertias, exactly as in (B.11d.1). Its cyclic index-$\ell$ subgroups
and the integral $p$-adic abelianizations used in B.11f are therefore
data of finite covers and quotients of $\Pi_U$, with their inherited
arithmetic actions. Finite extensions and good reduction may again
be compared with the completion by the same Galois comparison and
Tamagawa criterion as in B.11f. Alternatively all constructions below
can be made over the original henselian trait.

Suppose good marked reduction holds. The finite-cover construction
in B.11d, up to the application of AIK, uses only henselian lifting
and the prime-to-$p$ smooth-pair and Kummer comparisons, so it applies
to the present residue field. Explicitly, a smooth nonempty special
open has a point over a finite separable residue extension: use an
etale chart to affine space over the separable closure, then descend
the point. Henselian lifting and an unramified extension provide an
integral point. If there are at least two cusps, the modulo-$\ell$
abelian boundary quotient is nonzero at every cusp. If there is only
one cusp, positive genus and the Heisenberg $\ell$-group quotient
make its commutator nonzero. A characteristic refinement remains a
finite $\ell$-group quotient. The actual action based at the integral
point is unramified on this quotient; an additional finite unramified
extension kills its finite action and gives a constant-deck Galois
cover with a point above the chosen point. The genus estimate
(B.11d.2), with $\ell\ge5$, makes its compactification have genus
at least two. The regular tame charts $z=w^{e_c}$ give a good
proper model of that compactification. Each of these steps uses the
same SGA~1 statements and proofs specified in B.11d and does not use
residue perfection. B.11f now supplies $\mathcal P_\ell$.

Conversely, B.11f gives good reduction of the proper covering curve
$C$. Its constant deck group $Q$ extends faithfully to the special
fibre of its smooth stable model, by the finite unramified Isom
scheme of \citet[Theorem 1.11]{BC_DM}. Its order is invertible on
the trait. We spell out the one local issue needed over an imperfect
residue field in the tame quotient argument of B.11d.

At a geometric fixed point the stabilizer $H$ embeds faithfully in
the multiplicative group of the tangent line. A prime-to-$p$
finite-order automorphism with derivative one would have its first
nonzero higher coefficient multiplied by its order under iteration;
it must therefore be the identity on the formal local curve, hence
on the connected special curve. Faithfulness excludes this. Thus
$H$ is cyclic; if nontrivial its invariant tangent subspace is zero.
The fixed locus of $H$ is formally smooth over $R_L$ near this point.
Indeed, for a square-zero extension of test rings, smoothness first
gives a lift of a fixed point. The differences of its $H$-translates
form a cocycle in the relative tangent module; averaging, since
$|H|$ is invertible, changes the lift to an $H$-fixed lift. This
establishes the infinitesimal lifting criterion. The zero invariant
tangent space makes the fixed locus etale there. In particular its
residue extension is separable, so the required fixed section is
available after an etale base change even when the original residue
field is imperfect.

At that section the eigenparameter construction of B.11d gives
$h(w)=\zeta w$ and
\[
                       R_L[[w]]^H=R_L[[w^{|H|}]].
\]
The same computation with the semilocal product over an orbit gives
the local ring of the projective quotient $\mathcal C/Q$.
Thus this quotient is smooth and its reduced branch locus consists
etale-locally of disjoint sections. Every special fixed point lifts,
so no branch point occurs only in the special fibre. To check that
different branch sections cannot collide, translate their ramification
points by the deck group so that any alleged collision occurs at the
same special point of $\mathcal C$. Its cyclic stabilizer contains
both nontrivial generic stabilizers, since their fixed equations
persist under specialization. In the eigenparameter $w$, every
nontrivial subgroup of this tame stabilizer has the same unique
local fixed section $w=0$. The two purported distinct ramification
sections must therefore coincide; so do their images. Properness makes
the reduced branch divisor finite etale over $R_L$.
Its generic fibre is exactly $D_L$, since the cover is etale over
$U_L$ and ramified at every cusp. The quotient is therefore a good
marked model of $(\bar U_L,D_L)$.

Finally $L/K$ is unramified, so the identified inertia groups are
equal. Tamagawa's prime-to-$p$ criterion descends good marked
reduction to $R$. If one made the constructions after completion,
the same comparison used at the start descends the conclusion to
the original henselian trait. This proves the criterion in the
full stated scope.
\end{proof}


\subsection{B8--B10: effective actions and the contact map}

\begin{claim}{B.13}{Proved}
Let $(\mathcal X,\mathcal D)/R$ be a smooth marked model, let $x\in U(K)$,
and suppose its proper extension meets a rational special cusp $c$.
Write $z(x)=\varpi^\mu u$, where $z$ is a boundary parameter, $u\in R^*$
and $\mu>0$. The map
$\lambda:G_K\to\pi_1(\mathcal X\setminus\mathcal D)
\simeq\pi_1(U_k)$ induced by $x$ has inertia map
$T\xrightarrow{[\mu]}I_c\hookrightarrow\Pi$ and image of index $\mu$
in the cusp normalizer. Its restriction along $s_a(g)=(a(g),g)$ is the
cuspidal section of Kummer class
$\mu[a]+\kappa_k(\bar u)\in H^1(k,T)$.
These are statements after specialization to the total smooth model.
They are not statements that the generic point section has index $\mu$
in a generic cusp decomposition group.
\end{claim}
\begin{proof}
The prime-to-residue-characteristic Kummer covers of the punctured local
chart are dominated by $w^n=z$, by SGA 1, XIII, Propositions 5.1--5.2.
Pullback gives $w^n=\varpi^\mu u$. Units have $n$th roots over the strict
henselization, so inertia acts by $\mu$ times the uniformizer character.
The peripheral injection follows from the punctured-surface presentation,
including one cusp: its word is a product of commutators, whose nonzero
degree-two class in every free pro-$\ell$ quotient shows infinite order.
Thus $T\to\Pi$ is injective before abelianization.
The full map lands in $I_c\rtimes G_k$ and is
$(\tau,g)\mapsto(\mu\tau+b(g),g)$, where $[b]=\kappa_k(\bar u)$;
this follows by the same root computation on units. Its intersection with
$I_c$ is $\mu I_c$, and it surjects onto $G_k$, giving the index. Substituting
$\tau=a(g)$ proves the residue cocycle formula. A generic point section of
$E_{U_K}\to G_K$, by definition, has trivial intersection with its geometric
kernel. This proves that transferring the index assertion to that extension
would be a type error. In positive residue characteristic the same proof
uses $T_{p'}$, kills $P_K$, and replaces the index by the prime-to-$p$ part
of $\mu$.
\end{proof}

\begin{claim}{B.14}{Proved}
Suppose the residue outer action is realized by $V_s$ in B.4. A section
from a $K$-point or a rational cusp of $U$ restricts, after an explicitly
chosen finite residue extension, to a point or cusp action of $V_s$.
For an interior point meeting the special boundary the residue Kummer class
is the class in B.13; a point reducing away from the boundary gives its
ordinary special-fibre point section. A cuspidal section remains cuspidal,
with its Kummer class restricted along the chosen splitting.
\end{claim}
\begin{proof}
Choose a finite extension giving good marked reduction and then take
$H=s^{-1}(G_{K_1})$ and $G_L$ as in B.4. The section of the specified
point, or the specified cuspidal section, restricts through the Cartesian
arithmetic extension over $G_L$. Properness extends the point to the good
compactification. If its reduction avoids the boundary, functoriality and
smooth specialization identify its action with that of this special point.
If it meets the boundary, B.13 identifies the action with a cuspidal
section, including the unit and multiplicity contributions. A rational
generic cusp extends to a section of the etale boundary, and its normalizer
specializes to that cusp normalizer; the root calculation transports an
arbitrary cuspidal Kummer twist. Over a further finite residue extension the
descent datum defining $V_s$ is trivial, so these are literally point or cusp
actions on the base change of $V_s$. This is the exact sense of the remark
on cote 149, manuscript p.~51, scans 103--104: these constructions do
not produce a nongeometric genuine action above a geometrically realized
outer action.
\end{proof}

\begin{claim}{B.15}{Proved}
For a $K$-point specializing to a node $e$ of a split stable marked model,
$\Pi_e\subset C_\Pi((s_x\circ s)(G_k^\circ))$ for every residue splitting.
If the point specializes to a marked point, the corresponding cusp inertia
is contained in this centralizer.
\end{claim}
\begin{proof}
At a node the section factors through the local log point group of B.3.
If the point has local coordinates $x=\varpi^a u$,
$y=\varpi^{m-a}v$, its inertia coordinates are $(a\tau,(m-a)\tau,\tau)$,
with the unit cocycles on the arithmetic coordinates. On $G_k^\circ$ the
cyclotomic character is one. The whole rank-two Kummer kernel is abelian,
so all these elements commute with its antidiagonal $\Pi_e$, regardless
of $a$, $m$ or the splitting cocycle. At a cusp B.13 gives the same argument
with its rank-one kernel. The upper bounds are proved separately in
B.15u, using arithmetic cusp uniqueness and the profinite Bass--Serre tree.
In particular this local calculation does not identify outer fixed
conjugacy classes with elements fixed by the genuine action.
\end{proof}

\begin{claim}{B.15a}{Proved}
Put $F=k(\mu_\infty)$, where $k/\Q$ is finitely generated. For a
hyperbolic curve $Y/F$, an arbitrary continuous section of
$\pi_1(Y)\to G_F$ fixes at most one lifted geometric cusp. No extension
of this section to $G_k$ is required.
\end{claim}
\begin{proof}
We give the arithmetic input and the finite-cover argument explicitly.
For every finite extension $F'/F$ and every abelian variety $A/F'$, one has
\[
 T(A)^{G_{F'}}=0,\qquad \bigcap_{n\ge1}nA(F')=0.
\]
Here is the spreading argument extending Ribet's number-field theorem
\citep[Appendix, Theorem 1]{BC_KatzLang}. Write
$F'=k_0(\mu_\infty)$ with $k_0/\Q$ finitely generated and descend the
finite data to $k_0$. Let $a$ be its algebraic constant number field, and
spread $A$ to an abelian scheme over a smooth geometrically integral
$a$-variety $S$ with function field $k_0$. A rational torsion point over
$F'$ extends over $S_{a(\mu_\infty)}$: its finite etale division scheme
has a generic section, and normality extends that section. Evaluation at
a closed point of $S$, after passage to its number-field residue and its
cyclotomic closure, is injective on torsion sections: the equalizer of two
sections of a finite etale scheme is open and closed. Ribet's finiteness
therefore gives the first assertion. This argument also works after every
finite extension of $F$.

For the second assertion first consider a number-field cyclotomic closure
$a(\mu_\infty)$. A point divisible by every integer has compatible
factorial division points, by compactness of their finite nonempty sets.
Descend the abelian variety and the point to a finite subextension $a_1$.
The Kummer cocycle on $G_{a_1}$ vanishes on the normal subgroup
$G_{a(\mu_\infty)}$. Its values are therefore fixed by that subgroup,
so the first assertion makes the cocycle zero. All division points are
then defined over $a_1$. Reduction in a completion at a good finite place,
modulo successive powers of its maximal ideal, puts the point in the
intersection of the kernels of finite quotient groups. It is zero.
For general $F'$, spread a proposed nonzero infinitely divisible point as
a section over $S$ and remove its intersection with the identity section.
Every division point extends over $S_{a(\mu_\infty)}$ by the same finite
etale argument. Specialization gives a nonzero infinitely divisible point
over a number-field cyclotomic closure, a contradiction.

Suppose now that two distinct lifted cusps are fixed by the section
$\sigma:G_F\to\pi_1(Y)$. A finite geometric etale cover separates their
images. Refine it by a characteristic cover whose proper compactification
has genus at least two. Such a geometric cover exists: for a punctured
genus-zero curve, choose three punctures and restrict a Fermat quartic
cover of the resulting three-punctured line;
for genus one, an etale cover first supplies two boundary points and
a double cover branched there has genus two; in proper genus at least two
the identity suffices. A characteristic refinement of the Galois closure
preserves this genus bound by Riemann--Hurwitz. The subgroup generated by
this characteristic geometric subgroup and $\sigma(G_F)$ descends the
cover to $F$. The two fixed cusps become distinct $F$-rational points
$a,b$ of its proper compactification $\overline Y$.

Both points induce the same section of $\pi_1(\overline Y)\to G_F$,
because cusp inertia is killed on compactification. Under the Abel map
$z\mapsto[z-b]$, their difference has zero Kummer class for every
multiplication cover of $\Jac(\overline Y)$, hence
$[a-b]\in\bigcap_n n\Jac(\overline Y)(F)=0$. A function with divisor
$a-b$ would give a degree-one map to $\mathbb P^1$, contrary to the positive
genus. This proves uniqueness, including distinct lifts of the same base
cusp. It is the finite-neighbourhood proof of Stix's packet theorem,
with its arithmetic hypothesis verified over $F$
\citep[Theorem 14 and its proof; Appendix]{CF_Stix2008}.
\end{proof}

\begin{claim}{B.15u}{Proved}
Let $(\mathcal C,\mathcal D)/R$ be a split stable marked model in the
characteristic-zero setting of Part B. Let $x\in U(K)$ and let
$s:G_k\to G_K$ be any residue splitting. For the genuine action specified
by the point section, put
$H=(s_x\circ s)(G_k^\circ)$. Then
\[
 C_\Pi(H)=
 \begin{cases}
  1,&\text{if $x$ specializes to a smooth unmarked point},\\
  I_c,&\text{if $x$ specializes to a marked point $c$},\\
  \Pi_e,&\text{if $x$ specializes to a node $e$}.
 \end{cases}
\]
The inertia groups are the actual lifted subgroups selected by the local
point path. The equality is independent of the positive contact
multiplicity at a mark and of the two positive contact orders
$a,m_e-a$ at a node. In particular no other lifted node group or cusp
group is contained in this centralizer.
\end{claim}
\begin{proof}
We specify both the arithmetic vertex groups and the profinite-tree
argument; passing to an abstract tree or assuming that arbitrary fixed
subgraphs are path connected would not suffice.

Remove the open edges (the marked cusps) from the geometric semi-graph of
B.3, retaining the vertex groups and the closed-edge groups. This finite
graph of profinite groups has injective vertex and edge maps. Its standard
Bass--Serre graph $\mathcal T$ is a profinite tree, with
\[
 V(\mathcal T)=\coprod_v\Pi/\Pi_v,
 \qquad E(\mathcal T)=\coprod_e\Pi/\Pi_e.
\]
The sums are finite, so the edge set is compact. The geometric stabilizers
are the indicated conjugates. These are the construction and exact tree
theorem of \citet[\S3.7 and Proposition 3.8, pp.~419--420]{BC_ZM1989}.
The one-vertex graph with no nodes is included directly. The split
arithmetic action acts continuously on this tree and preserves the two
branches of each lifted edge; in particular there are no edge inversions.

For a lifted vertex $v$, let $Y_v$ be its normalized component minus the
nodes and marks. Its vertical logarithmic chart is the constant chart
$\mathbb N\to\mathcal O_{Y_v}$, $1\mapsto0$, pulled back from the base
log point. The Kummer covers of this chart are the roots of the base
parameter. After adjoining these roots the remaining finite cover is
etale on $Y_v$. Thus the arithmetic logarithmic vertex group $D_v$ has
the Cartesian description
\[
 D_v\simeq\pi_1(Y_v)\times_{G_k}G_K.
 \tag{B.vertex}
\]
The projection to $\pi_1(Y_v)$ kills precisely the vertical inertia,
which intersects the geometric vertex group trivially and centralizes it.
This is also the vertex part of the normalized-component gluing in B.3:
the chart is strict over the base log point away from its deleted special
points, and adjoining a root of the base chart leaves an ordinary finite
etale component cover. Conversely these component covers and the compatible
base roots give every such local Kummer cover. The injections in B.3
identify its geometric kernel with $\Pi_v$, and its projection onto $G_K$
is surjective. Its image is the entire arithmetic stabilizer: any other
stabilizer element can be multiplied by a local element with the same
image in $G_K$ to lie in $N_\Pi(\Pi_v)=\Pi_v$, the terminality assertion
of \citet[Proposition 1.2]{BC_MochComb}. Conjugating the chosen lift
conjugates this entire Cartesian description.

Pulling (B.vertex) back along $s$ identifies
$D_v\cap E_s$ with $\pi_1(Y_v)$. Consequently, for \emph{every}
$H$-fixed vertex, even one not fixed by the whole point-section image of
$G_k$, projection gives an actual continuous section
\[
 G_{k(\mu_\infty)}\longrightarrow
 \pi_1((Y_v)_{k(\mu_\infty)}).
\]
It projects to the identity on $G_{k(\mu_\infty)}$ because $H$ already
does so. The incident edges at $v$ are the lifted cusps of $Y_v$
corresponding to node branches; the identification follows from their
coset descriptions $\Pi_v/\Pi_e$. Distinct incident edges give distinct
lifted cusps, including when their images are the same branch at a finite
level. This identification is equivariant for $D_v$. Therefore two
distinct incident $H$-fixed edges would give two fixed lifted cusps for
the displayed section, contrary to B.15a. Every vertex of
$\mathcal T^H$ has degree at most one.

Theorem 2.8 of \citet[pp.~413--414]{BC_ZM1989} applies to an arbitrary
profinite group acting on a full profinite tree; it gives that
$\mathcal T^H$ is a subtree. It is nonempty: the local point path fixes
the vertex of the specialization if it is smooth, and fixes the lifted
node edge if the specialization is nodal. Its edge set is closed in the
compact edge set of $\mathcal T$, hence compact. A connected profinite
graph with compact edges and more than one vertex has an edge incident to
every vertex \citep[Lemma 2.14, p.~417]{BC_ZM1989}.
For clarity, this last fact follows by taking finite graph quotients:
the closed subsets of edges whose endpoint agrees with a given vertex
in each quotient have the finite intersection property.

It follows that $\mathcal T^H$ is either one vertex or one edge. Indeed,
in the nontrivial case every vertex has degree exactly one. The two endpoint
maps identify its vertex set with two disjoint copies of its compact edge
space: they are continuous bijections onto disjoint closed sets covering
all vertices, hence homeomorphisms. If the edge space had two points,
a nontrivial clopen partition would give a disconnected finite quotient
graph. This contradicts connectedness. This argument does not collapse
an edge and does not presume compactness of the edges after such a collapse.

Write $C=C_\Pi(H)$. It preserves $\mathcal T^H$, since its elements
commute with $H$. At a node the fixed tree already contains the local
edge, so it is exactly that edge. The action has no inversions, hence
$C$ lies in its geometric stabilizer $\Pi_e$. B.15 gives the reverse
inclusion for every splitting and for all contact orders.

If the specialization is smooth and unmarked, the action at its fixed
vertex, under (B.vertex), is the ordinary point action of $Y_v$ at that
special point. R15 gives $C\cap\Pi_v=1$. An incident fixed edge would
have its nontrivial inertia centralized by $H$: the edge action is the
cyclotomic character, which is trivial on $G_k^\circ$. Thus no such edge
exists, and the fixed tree is just $v$. Since $C$ preserves it,
$C\subset\Pi_v$ and $C=1$.

If the specialization is the mark $c$, the action at the fixed vertex
is the cuspidal action computed in B.13. R15 gives
$C\cap\Pi_v=I_c$. An incident fixed node edge would instead give a second
fixed lifted cusp of $Y_v$, contradicting B.15a. Again the fixed tree is
just $v$, and consequently $C=I_c$. Distinct lifted peripheral groups
are distinct maximal procyclic groups by the peripheral terminality in
\citet[Proposition 1.2]{BC_MochComb}; hence the displayed equalities
exclude every other node or cusp group. The multiplicities affect the
local Kummer homomorphism, but B.15 and the tree upper bounds show that
they do not affect this centralizer.
\end{proof}


\subsection{B11: reconstruction from inertia and arithmetic monodromy}

\begin{claim}{B.16}{Known}
Among stable marked models in this residue-characteristic-zero setting,
the dual semi-graph is functorially determined by the geometric group with
its $G_K$ outer action and its specified inertia subgroup. Once the
edge classes and Dehn coordinates are identified, the closed full-profinite
inertia image determines every positive integral thickness.
For proper curves the inertia image alone already determines the graph,
the full edge classes, and the thicknesses. For affine curves the same
conclusion holds if the geometric cusp classes are supplied as additional
data.
\end{claim}
\begin{proof}
First consider two proper stable curves and a continuous isomorphism of
their full geometric groups carrying one closed outer inertia image onto
the other. If those images are nontrivial, B.5 identifies each with $T$.
The isomorphism between them preserves their unique pro-$\ell$ Sylow
subgroups. In Example 2.5 of \citet{BC_MochComb} take the closed subgroup
$H=I_k^{\log}$ and its image $J$; this is allowed without requiring $H$
to be open. The action is $\ell$-logarithmically full because its
$\ell$-inertia is infinite. Both graphs are noncuspidal. Consequently
Corollary 2.7(iii), with $\Sigma$ the set of all primes, says that the
given isomorphism is graphic. This recovers the full edge classes, not
merely their images in one pro-$\ell$ quotient. If the inertia image is
trivial, Tamagawa implies that the stable model is smooth, whose one-vertex
graph is determined by the geometric genus. For affine curves with cusp
classes given, apply instead \citet[Corollary 4.2]{BC_HMNodes} to the
inertia representations: they have NN-type by B.5 and the positive
thicknesses, and the isomorphism is group-theoretically cuspidal by
hypothesis. That corollary gives the same full graphic conclusion.

We also check the arithmetic reconstruction route specified in L8.
Choose a prime $\ell$ and pass through the characteristic quotient
$\Pi\twoheadrightarrow\Pi^{(\ell)}$. An equivariant isomorphism of full
groups induces an equivariant isomorphism of these quotients and of their
outer image groups. In the nodal case B.5 gives infinite $\ell$-inertia,
which is the logarithmic fullness required in
\citet[Corollary 2.8]{BC_MochComb}. The action is arithmetically full because
it is the image of the whole log Galois group. The residue field embeds in
a finitely generated extension of $\Q_r$ for a prime $r\ne\ell$: choose a
transcendence basis over a number field, embed that number field in a finite
extension of $\Q_r$, and take a finite extension containing the remaining
algebraic generators. Thus the prime-set hypothesis of that corollary is
satisfied with $\Sigma=\{\ell\}$, not incorrectly with all primes.
The given compatibility with the specified $I_K$ matches $\ell$-inertia.
The corollary says the isomorphism is graphic and hence reconstructs the
dual semi-graph. If inertia is trivial the model is smooth by Tamagawa;
the single vertex genus and cusp number are recovered from homology,
properness and Stix's cusp criterion as in B.4.

With a graphic identification in hand, Theorem 4.8(iv) of
\citet{BC_HMTopics} identifies the Dehn group with $\prod_eT$.
Projection of the inertia image to the $e$ coordinate is $m_e\widehat\Z$.
This ideal determines all $v_\ell(m_e)$, hence the unique positive integer
$m_e$. A chosen generator can change by a unit of $\widehat\Z$ without
changing those ideals. If only prime-to-$p$ coordinates were retained,
the $p$-part of $m_e$ would be lost. This proves the precise conditional
thickness assertion. The affine extraction without cusp data is proved
next by a different reconstruction theorem.
\end{proof}

\begin{claim}{B.16u}{Proved}
Let $(\mathcal C,\mathcal D)/R$ be a stable marked model in the
residue-characteristic-zero setting of Part B. Put
$J=\rho(I_K)\subseteq\Out(\Pi)$, with the full geometric profinite
group $\Pi$. The pair $(\Pi,J)$ determines the set
\[
 \{([\Pi_e]_{\Pi},m_e):e\text{ a geometric node}\}.
\]
This holds for affine curves without a supplied set of cusp classes,
and is functorial under every continuous isomorphism of the geometric
groups carrying one closed inertia image onto the other. No generator
of $J$ is part of the data. The thicknesses use the normalized valuation
of the trait over which the stable model is given.
\end{claim}
\begin{proof}
An unramified finite extension splitting the special fibre changes neither
geometric groups, geometric node classes, thicknesses nor inertia. We can
therefore split the fibre for the proof without adding data to the
reconstruction. If $J=1$, B.5 says there are no nodes. Suppose there is a
node. Then B.5 identifies $J$ with $\widehat\Z$ and, after a complex
topological comparison and a choice of orientation, identifies a
topological generator with the completion of
\[
 \varphi=\prod_e\tau_e^{m_e}\in\operatorname{Mod}(S),
\]
where $S$ is a compact oriented surface whose boundary circles represent
the cusps. The node circles form a disjoint essential multicurve, and
every component $S_v$ cut out by that multicurve has $\chi(S_v)<0$.
These facts are exactly stability: a disk or an annulus component would
be an unstable rational component, and a peripheral or parallel pair
of node circles would cut off such a component. A base point outside
the twist annuli gives an actual automorphism of
$\Gamma=\pi_1(S)$. We use this choice only to prove the reconstruction
formula; it is not retained as input.

\paragraph{The group constructed from the inertia image.}
Because $\Pi$ is finitely generated and centre-free, B.2 constructs,
from the given pair alone, the profinite extension
\[
 1\longrightarrow\Pi\longrightarrow
 E_J:=\Aut(\Pi)\times_{\Out(\Pi)}J
 \stackrel{q}{\longrightarrow}J\longrightarrow1.
\]
The chosen automorphism $\widehat\varphi$ generates a procyclic
subgroup of $\Aut(\Pi)$ which maps isomorphically onto $J$:
the composite $\widehat\Z\to\Aut(\Pi)\to J$ is an isomorphism
by B.5. Consequently
\[
 E_J\simeq\Pi\rtimes_{\widehat\varphi}\widehat\Z
       \simeq\widehat{\Gamma\rtimes_{\varphi}\Z}.
\]
For completeness, the second isomorphism is full profinite completion,
not a procyclic approximation. Characteristic finite-index subgroups
of the finitely generated group $\Gamma$ are cofinal among its
finite-index subgroups. Thus $\varphi$ extends to the indicated
continuous action on $\widehat\Gamma=\Pi$. Every homomorphism
$\Gamma\rtimes\Z\to Q$ to a finite group extends on $\Gamma$ to
$\Pi$ and on $\Z$ to $\widehat\Z$; the semidirect relation extends
by continuity. Conversely their product has dense image. These facts
give both universal maps and show they are inverse. The marked normal
subgroup $\Pi\subset E_J$ is preserved throughout.

\paragraph{The mapping torus and its geometric decomposition.}
Let $M_\varphi$ be the mapping torus of the multitwist, represented by a
homeomorphism equal to the identity off the disjoint twist annuli and on
$\partial S$. It is compact and orientable, with one boundary torus
for each cusp, and its fundamental group is
$\Gamma\rtimes_\varphi\Z$. The interior has universal cover
$\widetilde{\operatorname{int}S}\times\R\simeq\R^3$.
Irreducibility descends from this cover by
\citet[Proposition 1.6]{BC_Hatcher3M}. Each boundary torus is
incompressible: in the surface-bundle extension its group is generated
by a nontrivial boundary loop and a lift of the generator of $\Z$,
and projection to $\Z$ first kills the second exponent in any proposed
relation, after which surface-group injectivity kills the first.
The same argument applies to each node torus.

Cutting at the node tori gives the product Seifert pieces
$M_v=S_v\times S^1$. There are no solid-torus, annular-product or
twisted-Klein-bottle pieces since $\chi(S_v)<0$. These tori are exactly
the JSJ tori. Here are the relevant minimality checks. None is
boundary-parallel, and no two are parallel: either possibility would
make the corresponding essential surface circles peripheral or
freely homotopic, respectively. Such circles are excluded by the stable
multicurve. On the torus associated to $e$, choose the surface circle
$c_e$ and the circle fibre $t_v$ on one side. The circle fibre on the
other side is
\[
             t_w=c_e^{\epsilon_e m_e}t_v,
             \qquad \epsilon_e\in\{1,-1\}.
\tag{B.JSJ.1}
\]
Indeed the two product trivialisations of the annulus mapping torus
differ by exactly $m_e$ full twists. This also follows directly from
the separating-edge or HNN formula in B.5. Thus the two primitive
fibre slopes have intersection number $m_e>0$.

Deleting any of these cuts leaves an essential torus, so the joined
piece is not atoroidal. It cannot be Seifert either. In an irreducible
Seifert piece an essential torus is vertical or horizontal
\citep[Proposition 1.12]{BC_Hatcher3M}. A horizontal torus gives a
finite torus cover of the base orbifold and hence a virtually solvable
fundamental group, whereas the joined piece contains the nonabelian
free group $\pi_1(S_v)$. It must therefore be vertical. The induced
Seifert fibrations on each $S_v\times S^1$ have the product fibre
slope on the boundary: the only exceptions to boundary-slope uniqueness
are a solid torus or an $I$-bundle over a torus or Klein bottle
\citep[Lemma 1.17]{BC_Hatcher3M}, already excluded here.
Their slopes would have to coincide across the omitted cut, contrary
to (B.JSJ.1). This argument also applies to several omitted cuts and
to a loop edge. Minimality and uniqueness in the torus-decomposition
theorem identify the displayed decomposition with the JSJ decomposition
\citep[Theorem 1.9]{BC_Hatcher3M}.

\paragraph{The exact profinite JSJ input.}
We use \citet[Theorem 5.5]{BC_Xu2025}, in the opened arXiv version~4,
whose hypotheses are compact orientable irreducible $3$-manifolds with
empty or incompressible toral boundary, excluding closed Sol manifolds.
It states that every isomorphism of their profinite fundamental groups
induces a congruent isomorphism of their JSJ graphs of profinite groups.
In particular it preserves vertex and edge groups up to ambient
conjugacy, with the two incidence homomorphisms. No peripheral marking
is a hypothesis. The same source, Propositions 5.2--5.4, supplies
efficiency, the standard profinite tree and the boundary version used
in that theorem. Our manifold satisfies these hypotheses; when there
are nodes its nonabelian product pieces exclude the Sol exception.
This is a different input from the cuspidal or noncuspidal hypotheses
of the combinatorial anabelian results used in B.16.

Accordingly $E_J$ recovers the profinite groups
\[
 G_v=\widehat{\pi_1(S_v)}\times\widehat\Z,
 \qquad A_e=\widehat{\pi_1(T_e)}\simeq\widehat\Z^2,
\]
and their incidences. The geometric node group is recovered by
\[
                         \Pi_e=A_e\cap\Pi.
\tag{B.JSJ.2}
\]
This is an equality with the full closed cyclic subgroup: the surface
bundle map on $A_e=\overline{\langle c_e,t_v\rangle}$ sends $c_e$
to zero and $t_v$ to a generator of $J$, so its kernel is precisely
$\overline{\langle c_e\rangle}$.

The JSJ theorem initially gives ambient $E_J$ conjugacy. This yields
exactly the required $\Pi$ conjugacy. For any $g\in E_J$, choose
$h\in A_e$ with $q(h)=q(g)$; such an $h$ exists because
$q(A_e)=J$. Then $gh^{-1}\in\Pi$, and $h$ normalizes
$A_e\cap\Pi$. Hence the $E_J$ conjugates of this node group are
already its $\Pi$ conjugates. The same adjustment applies after an
isomorphism between two reconstructed extensions.

\paragraph{An intrinsic integral thickness formula.}
The surface factor of $G_v$ is a free profinite group of rank at least
two, hence centre-free; thus $Z(G_v)$ is exactly its product
$\widehat\Z$ fibre. For the two incidence embeddings
$\iota_{e,i}:A_e\hookrightarrow G_{v_i}$, put
\[
 F_{e,i}=\iota_{e,i}^{-1}(Z(G_{v_i})),\qquad i=0,1.
\]
Both are closed primitive cyclic subgroups of $A_e$. They remain two
different subgroups when $v_0=v_1$, because the two incidence embeddings
are retained. Formula (B.JSJ.1), completed profinitely, gives
\[
 \boxed{\quad m_e=[A_e:F_{e,0}F_{e,1}].\quad}
\tag{B.JSJ.3}
\]
In the basis $(c_e,t_v)$ the two columns are $(0,1)$ and
$(\epsilon_em_e,1)$, so their cokernel is
$\widehat\Z/m_e\widehat\Z$, of order $m_e$.
This finite index is preserved by every congruent isomorphism of the
JSJ groups. It uses no orientation or chosen generator of $J$.
It also explains directly why retaining all primes recovers the full
positive integer rather than only selected valuations.

Finally an isomorphism $\alpha:\Pi\to\Pi'$ carrying $J$ onto $J'$
induces $E_J\to E_{J'}$ by conjugation on automorphisms. Applying the
JSJ theorem, (B.JSJ.2) and (B.JSJ.3) proves the stated functorial
reconstruction. If a ramified extension of degree $d$ is specified,
its inertia image is the corresponding open subgroup and the same
construction gives the base-changed thickness $dm_e$, as it must.
No thickness independent of the normalization of the trait is asserted.
\end{proof}


\subsection{B12--B14: the three explicit tests}

\begin{claim}{B.17}{Proved}
For the Tate elliptic curve $E_q$ over $\Q((q))$ with its origin removed,
the stable marked special fibre is a rational nodal curve with one mark,
one loop edge of thickness one and one cusp. In the presentation
$\Pi=\widehat{\langle a,b\rangle}$, node inertia is $\overline{\langle a\rangle}$,
cusp inertia is generated by $[a,b]^{-1}$, and positive inertia is
$a\mapsto a$, $b\mapsto ba$. The graded dimensions of B2 are $(1,0,1)$
in weights $(0,-1,-2)$. All five conditions $a$--$e$ are false.
\end{claim}
\begin{proof}
The Tate plumbing identifies the two branches of the normalization by
$xy=q$; its special normalization is $\mathbf P^1$ with branch points
$0,\infty$ and the origin mark $1$. It is stable, has one node and the
stated thickness. Cutting the torus along $a$ gives the HNN presentation;
the single crossing $b$ gives the stated twist. Its matrix on homology in
the basis $(a,b)$ is $\left(\begin{smallmatrix}1&1\\0&1\end{smallmatrix}\right)$.
Thus $b,c,d$ fail and so does $a$. The weight-zero quotient is nonzero, so
B.4 excludes $e$ for every splitting. More explicitly, the Tate-module
extension is $0\to\Q_\ell(1)\to V_\ell E_q\to\Q_\ell\to0$ with Kummer
class of $q$. On the reference splitting fixing its roots this class is
zero; on $s_a$ it is the $\ell$-component of $[a]$. This verifies explicitly
both the possible change of the extension and independence of its graded
pieces. The surface matrix calculation is checked in the code.
\end{proof}

\begin{claim}{B.18}{Proved}
For $U:y^2=x^5-t$ minus infinity over $\Q((t))$, good marked reduction is
obtained over $r^{10}=t$ by $x=r^2X$, $y=r^5Y$. The special pair is
$Y^2=X^5-1$ with its unique infinity, so there are no nodes after this
extension. Outer inertia has order ten and acts on homology with
characteristic polynomial $\Phi_{10}(z)=z^4-z^3+z^2-z+1$.
The answers to $a,b,c,d,e$ are respectively no, no, no, no, yes.
For $[s-s_0]\mapsto a\in\Q^*/\Q^{*10}$, the realizing curve is
$y^2=x^5-a$ minus infinity, with the cocycle convention of B.9.
\end{claim}
\begin{proof}
The indicated substitution gives $Y^2=X^5-1$ exactly; the degree-five
polynomial is square-free in characteristic zero, and its smooth proper
completion has genus two and a single rational point at infinity.
A generator $r\mapsto\zeta_{10}r$ acts on the special coordinates by
$(X,Y)\mapsto(\zeta_{10}^{-2}X,\zeta_{10}^{-5}Y)$, of order ten.
On holomorphic differentials $dX/Y$, $X\,dX/Y$ the eigenvalues are
$\zeta_{10}^3,\zeta_{10}$, and on full homology the eigenvalues are the
four primitive tenth roots. Thus the order is already ten on homology,
and $b,c,d$ fail. Tamagawa excludes $a$; B.4 gives $e$.
The sole graded piece after the good extension is weight minus one of
dimension four.

An explicit group description retaining the peripheral element is useful:
let $\Gamma=\widehat{C_5*C_2}=\langle u,v\mid u^5=v^2=1\rangle$ and map
it onto $C_{10}$ by $u\mapsto4$, $v\mapsto5$. Then
$\Pi=\ker(\Gamma\to C_{10})$ is free profinite of rank four, and its
cusp group is generated by $(uv)^{10}$, with inverse permitted by orientation.
This is the orbifold covering $X^5=z$, $Y^2=z-1$ of the affine line, with
orbifold orders five and two at $0,1$. The kernel is torsion-free and has
Euler characteristic $10(1/5+1/2-1)=-3$, proving its rank. Conjugation by
$u^4v$, whose image is $1\in C_{10}$, describes the positive inertia outer
class, up to the chosen orientation convention. Finally if $b^{10}=a$,
the invariant coordinates $x=b^2X$, $y=b^5Y$ for the twisted descent datum
satisfy $y^2=x^5-a$. This proves the twist formula for every splitting class.
The substitutions and cyclotomic polynomial are verified symbolically.
\end{proof}

\begin{claim}{B.19}{Proved}
For $\mathbf P^1\setminus\{x^3=t\}$ over $\Q((t))$, good marked reduction
is obtained over $r^3=t$, $x=rX$, with special pair
$(\mathbf P^1,\{X^3=1\})$. There are no nodes after this extension. The
group is free profinite of rank two with cusp generators
$c_0,c_1,c_2$ and relation $c_0c_1c_2=1$; inertia cyclically permutes the
three cusp conjugacy classes and has order three in the outer group.
Its homology polynomial is $z^2+z+1$. Conditions $a,b,c,d,e$ have values
no, no, no, no, yes. If $[s-s_0]$ maps to $a\in\Q^*/\Q^{*3}$, then
$V_s=\mathbf P^1\setminus\{x^3=a\}$.
\end{claim}
\begin{proof}
The substitution is exact, and the divisor $X^3=1$ is finite etale over
$\Q$, though its points split only over $\Q(\zeta_3)$. Inertia acts by
$X\mapsto\zeta_3^{-1}X$. It permutes the cusp classes cyclically; on their
augmentation module the nontrivial eigenvalues are the primitive cube
roots. Thus the homological action has order three. This proves failure of
$b,c,d$, and failure of $a$ follows either from the ramified divisor or
Tamagawa. Good reduction after the cubic extension gives $e$.
Over $\Q(\zeta_3)$ the B2 piece is $\Q_\ell(1)^2$; over $\Q$ it retains
the Galois permutation action on the three roots, as required by the
nonsplit form of B.6. Twisting $X$ by a cube root of $a$ gives precisely
the displayed divisor. This example checks that ``good reduction of the
curve'' cannot replace good reduction of the marked pair, even in type
$(0,3)$.
\end{proof}

\subsection{The full group in equal characteristic: completion of B3}

In this subsection $R$ is a henselian DVR of equal characteristic $p>0$,
$K=\operatorname{Frac}(R)$, and $k$ is its residue field. The conditions
$a,b,c,d,e$ have the meanings in B3, with the \emph{full} geometric
fundamental group in $b,c,e$, and $\ell\ne p$ in $d$. A residue splitting
is assumed given whenever $e$ is mentioned. These are modern extensions
of cote 149, manuscript pp.~49--51, scans 99--104; we do not attribute
the positive-characteristic conclusions below to the manuscript.

\begin{claim}{B.20}{Proved}
For a proper hyperbolic curve $C/K$, trivial full outer inertia implies
that $C$ has good reduction and that the full geometric specialization map
\[
 \operatorname{sp}:\pi_1(C_{\bar K})\longrightarrow
                  \pi_1(\mathcal C_{\bar k})
\]
is an isomorphism. Its outer equivariance identifies the action restricted
along \emph{every} residue splitting with that of the smooth special
fibre. In particular $b\Rightarrow e$.
For an affine hyperbolic curve in equal characteristic $p$, full inertia
acts nontrivially already on $H^1_{\mathrm{et}}(U_{\bar K},\F_p)$.
Consequently full $b$ and $c$ are impossible in the affine case.
\end{claim}
\begin{proof}
We first prove the proper assertion. Trivial full outer inertia implies
triviality on every characteristic pro-$\ell$ quotient. Tamagawa's
good-reduction criterion therefore gives a smooth proper model
$\mathcal C/R$ \citep[Theorem 5.3]{BC_Tamagawa}. The specialization map
is surjective: SGA 1, X, Corollary 2.4 applies to every proper smooth
morphism over a locally noetherian base, not just to a complete trait
\citep[X, Corollary 2.4]{BC_SGA1}. We prove its injectivity by treating
\emph{all} finite geometric Galois covers, including those of degree
divisible by $p$.

Let $D_{\bar K}\to C_{\bar K}$ be a connected finite etale Galois cover
with group $Q$ and degree $d$. Finite presentation allows us, after a
finite separable extension $L/K$, to descend the cover, all its deck
automorphisms and a point of $D$ to $L$. The latter gives compatible
sections of the arithmetic sequences of $D$ and $C$, and hence actual
Galois actions on $H=\pi_1(D_{\bar K})\triangleleft
\Pi=\pi_1(C_{\bar K})$. For $\sigma\in I_L$, the action on $\Pi$ is
inner by hypothesis. The induced outer action on $H$ consequently belongs
to the finite image of $\Pi/H$ in $\Out(H)$: conjugation by an element
of $H$ is inner on $H$. This conclusion does not require a choice of
continuous conjugating elements in $\Pi$.

The continuous outer action of $I_L$ on $H$ thus has finite image. Its
kernel is open in $I_L$; choose an open subgroup of $G_L$ whose
intersection with $I_L$ lies in that kernel. After replacing $L$ by the
corresponding finite extension, the outer inertia action on $H$, and
therefore on its pro-$\ell$ quotient, is trivial. Applying the same
good-reduction criterion to the proper curve $D$ gives a smooth proper
model $\mathcal D/R_L$. Here $g(D)=1+d(g(C)-1)\ge2$.

We now justify extension of the \emph{map}; good models alone would not
suffice. The chosen $L$-point of $D$ extends by properness and gives an
$R_L$-section of $\mathcal D$ and, by its image on the generic fibre,
an $R_L$-section of $\mathcal C$. The relative Abel map embeds
$\mathcal C$ as a closed subscheme of its relative Jacobian $J$.
The Neron mapping property extends the generic map
$\mathcal D_L\to\mathcal C_L\hookrightarrow J_L$ uniquely across
the smooth $R_L$-scheme $\mathcal D$.
For precise sources, SGA 7 I, IX, 1.1 defines that mapping property on
\emph{all} smooth schemes over the trait, and Theorem 12.1(c)--(d),
especially (12.1.8)--(12.1.10), identifies the relative degree-zero Picard
scheme with the identity component of the Neron model
\citep[IX, 1.1 and Theorem 12.1]{BC_SGA7}. One may apply the mapping
property to the Neron model first: the extended section maps to zero,
and the connected special fibre of $\mathcal D$ then maps into its
identity component $J$. These identifications can be checked over the
strict henselization and descend by their uniqueness.
The extended map factors through the closed subscheme $\mathcal C$:
the pulled-back defining ideal vanishes on the dense generic fibre and
$\mathcal O_{\mathcal D}$ is torsion-free over the DVR. We obtain
$f:\mathcal D\to\mathcal C$.

Choose a relatively ample line bundle $M$ on $\mathcal C$. The degree
of $f^*M$ on fibres of the proper flat family $\mathcal D/R_L$ is
constant, by constancy of its Euler characteristic. Its generic degree
is positive. Therefore the map on the special fibre is nonconstant.
All its fibres are finite, so $f$ is proper and quasi-finite, hence finite.
At a point of $\mathcal D$ above a point of $\mathcal C$, both local
rings are regular and have the same dimension. A system of parameters
from the latter is a system of parameters in the former because the
map is finite. Since a regular local ring is Cohen--Macaulay, this is
a regular sequence, and the local flatness criterion gives flatness.
Thus $f$ is finite flat of rank $d$.

Every generic deck automorphism extends uniquely to $\mathcal D$.
Here the exact input is Deligne--Mumford's Theorem 1.11: the isomorphism
scheme of two stable curves is finite and unramified; its properness
extends the automorphism over the DVR
\citep[Theorem 1.11, p.~84]{BC_DM}. Specialization of these automorphisms
is injective. Indeed the equalizer of two sections of an unramified
separated scheme is open and closed; if it contains the closed point
of a local trait, it is the entire trait. They commute with $f$ by
uniqueness on the dense generic fibre. The special map, of total degree
$d=|Q|$, consequently has $d$ distinct deck automorphisms. Its separable
degree is therefore at least $d$, so it is a separable Galois map of
degree $d$ and has no inseparable factor.

The genera of the smooth proper fibres are constant, and the generic
cover is etale. Hence
\[
 2g(\mathcal D_k)-2=d\bigl(2g(\mathcal C_k)-2\bigr).
\]
Riemann--Hurwitz on the special fibres makes its effective different
zero. Thus the special map is etale. A finite flat morphism which is
etale on both fibres of a trait is etale; this follows either from
the fibrewise criterion or from the vanishing of its module of
relative differentials. We have proved that the given geometric cover
extends etale after a finite extension of the trait.

Specialization is compatible with that extension and with this finite
etale cover. Accordingly the quotient $\Pi\to Q$ factors through
$\pi_1(\mathcal C_{\bar k})$. The change of algebraically closed fields
does not change these \emph{proper} geometric groups
\citep[X, Corollary 1.8]{BC_SGA1}. Every finite quotient of $\Pi$ has
now been treated; their kernels have trivial intersection, proving
injectivity of $\operatorname{sp}$. The construction of specialization
from the arithmetic fundamental-group diagrams in SGA 1, X, 2.1--2.4
also proves its outer $G_K$-equivariance, with $G_K$ acting on the
special group through $G_k$. Since $\operatorname{sp}$ is an isomorphism,
restriction along any $s:G_k\to G_K$ is exactly the special-fibre
outer action. This proves $e$ over $k$ itself.

For the affine assertion, passing to a finite separable extension
preserves any putative triviality of inertia. We may therefore split a
boundary point $P$ and choose, by Riemann--Roch, a function
$f\in\mathcal O(U)$ whose only pole is at $P$, of order $m>0$ prime
to $p$. We choose $m$ sufficiently large. If $\varpi$ is a uniformizer,
let $a$ satisfy $a^p-a=\varpi^{-1}$. This defines a totally ramified
Artin--Schreier extension of degree $p$: the valuation of $a$ is
$-1/p$, so the extension has ramification index $p$. Some
$\tau\in I_K$ satisfies $\tau(a)=a+1$. On the affine geometric curve,
the Artin--Schreier sequence gives
\[
 H^1_{\mathrm{et}}(U_{\bar K},\F_p)
 =\mathcal O(U_{\bar K})/(h^p-h).
\]
Then $\tau[af]-[af]=[f]\ne0$. To check the inequality, if
$h^p-h=f$, any pole of $h$ gives a pole of $f$ of order divisible by
$p$, and no pole can disappear. The unique pole of $f$ has order $m$
prime to $p$, a contradiction. Thus the inertia action on the mod-$p$
abelian quotient is nontrivial, proving the last assertion.
\end{proof}

\begin{claim}{B.21}{Proved}
Let $k=\F_9=\F_3(\alpha)$ with $\alpha^2=-1$, $R=k[[t]]$,
$K=k((t))$, and $\lambda=-1+t$. The smooth proper genus-two curve
\[
 C:\quad y^2=(x^2+\alpha)(x^2+\alpha-1)(x^2+\alpha-\lambda)
\]
has $a,d$ true and $b,c,e$ false. Failure of $e$ holds for every
residue splitting and after every finite residue extension.
\end{claim}
\begin{proof}
Consider the smooth elliptic model
\[
 E:\quad v^2=u(u-1)(u-\lambda)=u^3-tu^2+\lambda u.
\]
Its discriminant $16\lambda^2(1-\lambda)^2$ is a unit. For a
characteristic-three cubic $v^2=u^3+a_2u^2+a_4u+a_6$, the nonzero
three-torsion points are the flexes, with
$a_2u^3+a_2a_6-a_4^2=0$. This is also the explicit calculation of
\citet[Example 5.27]{BC_MoonenIsogs}. In our case their $u$ coordinates
satisfy $u^3=-\lambda^2/t$.

We use this formula on the \emph{etale Verschiebung kernel}, avoiding
the imperfection pitfall in substituting $K^{\mathrm{sep}}$ for $\bar K$
in $E[3]$. Let
$F:E\to E^{(3)}$ be relative Frobenius and
$V:E^{(3)}\to E$ Verschiebung, so $VF=[3]$; these maps and the
identity over a characteristic-$p$ base are constructed in
\citet[Definitions 5.18 and Propositions 5.19--5.20]{BC_MoonenIsogs}.
The two nonzero points of $\ker V$ have coordinates
\[
 X=-\frac{\lambda^2}{t},\qquad
 Y^2=X(X-1)(X-\lambda^3)
     =-\frac{\lambda^4(\lambda-1)^4}{t^3}.
\]
Indeed Frobenius is radicial and bijective on algebraically closed
points, and carries the two flexes to these points; $\deg V=3$.
On putting $z=Yt^2/(\lambda^2(\lambda-1)^2)$ we obtain
\[
 z^2=-t.
\]
This quadratic polynomial is separable and Eisenstein. Thus the two
nonzero points of the \emph{finite etale} $K$-group scheme $\ker V$
are exchanged by inertia. Notice that the assertion concerns a cover
$E^{(3)}\to E$ which is etale over $K$, not a possibly nonetale
$3$-torsion subgroup of $E$ over this imperfect field; the distinction
is explicitly warned about in \citet[Caution 5.24]{BC_MoonenIsogs}.

The map $C\to E$, $(x,y)\mapsto(x^2+\alpha,y)$, has degree two.
The displayed sextic expands over $\F_9[[t]]$ as
\[
 f(x)=x^6-tx^4+\bigl((\alpha+1)t-1\bigr)x^2
                  +\alpha+(\alpha+1)t.
\]
At $t=0$ its derivative is $x$, while $f(0)=\alpha\ne0$.
The two points at infinity are smooth since the leading coefficient
is one and $2$ is invertible. Hence the proper model is smooth of genus
two, proving $a$ and $d$.
Pulling back the degree-three etale Verschiebung cover along $C\to E$
gives a geometrically connected cover: the two function-field degrees
two and three are coprime, and the latter extension is Galois over
the algebraic closure. Thus the geometric abelianized group of $C$
has the inertia-ramified quotient $\ker V(\bar K)$.
It follows that $c$, and therefore $b$, is false.

We compute the geometric $3$-ranks to exclude $e$. On the holomorphic
basis $\eta_0=dx/y$, $\eta_1=x\,dx/y$, Cartier's formula gives
\[
 \mathcal C(\eta_0)=\bigl((\alpha+1)t-1\bigr)^{1/3}\eta_0,
 \qquad
 \mathcal C(\eta_1)=(-t)^{1/3}\eta_1.
\]
Here we compute over an algebraic closure: write $1/y=f(x)/y^3$
and use $\mathcal C(x^{3j+2}dx)=x^jdx$. The semilinear map is invertible
generically and has stable rank one at $t=0$. Thus
$f(C_{\bar K})=2$ and $f(C_{\bar k})=1$.

Suppose the full residue action were realized by a smooth hyperbolic
curve $W$ over a finite extension $k'/k$. The intrinsic pro-prime-to-$3$
quotient then has the same outer action as $C_{k'}$: good proper
specialization identifies it independently of the splitting. Cohomology
at any $\ell\ne3$ distinguishes a proper curve from an affine curve
and determines its genus; hence $W$ is proper of genus two. Centre-free
extension reconstruction B.2 identifies the associated geometrically
prime-to-$3$ arithmetic groups. Saidi--Tamagawa's exact finite-field
anabelian theorem applies: their Theorem 3.9 assumes all primes other
than the field characteristic in the geometric group and concludes that
an isomorphism of those arithmetic groups comes from a scheme
isomorphism \citep[Theorem 3.9]{A_SaidiTamagawa2009}.
Consequently $W$ is isomorphic to $C_{k'}$ as a scheme; a possible
automorphism of its finite field of constants does not affect its
$3$-rank. But its full geometric group was assumed isomorphic to
$\pi_1(C_{\bar K})$, and
$\dim_{\F_3}H^1(\pi_1(-),\F_3)$ is the $3$-rank. It would have
to be both two and one. This proves failure of $e$, for all the
splittings and finite residue extensions asserted.
\end{proof}

\begin{claim}{B.22}{Proved}
Over $\F_7((t))$ the proper genus-two curve
\[
 D:\quad y^2=(x^3-t^6)(x^3-1)
\]
has $c,d$ true and $a,b,e$ false for the full-group conditions.
\end{claim}
\begin{proof}
The normalized compact-type model is the same explicit model as in
B.11, with uniformizer $t$. We recall its checks. In the inner chart
$x=t^2X$, $y=t^3Y$, the special component is $Y^2=1-X^3$.
The outer component normalizes to $V^2=x(x^3-1)$, which becomes
$Y^2=1-X^3$ under $X=1/x$, $Y=V/x^2$. Near the connecting point put
$xz=t^2$ and $A=(1-z^3)(x^3-1)$. The normalization has
\[
 V^2=xA,\qquad W^2=z/A,\qquad VW=t.
\]
Since $A(0,0)=-1$, this is a regular separating node of thickness one.
Both components are smooth elliptic curves in characteristic seven,
and the node gives a stable compact-type special fibre. The model
therefore remains nodal after every finite extension, so $a$ is false
and Tamagawa makes $b$ false. C.5 excludes even prime-to-$7$ residue
realization, hence excludes full $e$.

The relative Jacobian $J$ is an abelian scheme with special fibre the
product of the two component Jacobians. Each is ordinary: the Hasse
coefficient of $x^6$ in $(1-x^3)^3$ is $3\ne0$ in $\F_7$.
For $\ell\ne7$ good reduction gives trivial inertia on $T_\ell J$.
For the residue prime we give the additional argument rather than
appealing to a prime-to-$p$ criterion. The iterated Verschiebung
\[
 V^n:J^{(7^n)}\longrightarrow J
\]
is finite flat, of rank $7^{2n}$, and its special kernel is etale because
the special abelian surface is ordinary. Thus its kernel is finite
etale over $R$: the nonetale locus of this finite flat scheme is closed,
its image in the local trait is closed, and it misses the closed point.
Over an algebraically closed field, relative Frobenius is a bijection
on points and identifies these kernels with the etale $7^n$-torsion
quotients. They consequently give unramified $7$-primary geometric
Tate modules. More directly, the etale isogenies $V^n$ form the
cofinal $7$-power covers of the ordinary abelian variety, all extended
over $R$. Hence inertia is trivial on every primary part of
$\pi_1(D_{\bar K})^{\mathrm{ab}}=\pi_1(J_{\bar K})$.
This proves $c$ and $d$.
\end{proof}

\begin{claim}{B.23}{Proved}
Over $K=\F_3((t))$ the proper genus-two curve
\[
 C:\quad y^2=t(x^5-x)
\]
has $e$ true for every residue splitting, whereas $a,b,c,d$ are false.
\end{claim}
\begin{proof}
The constant curve $C_0:Y^2=x^5-x$ over $\F_3$ is smooth of genus two:
the polynomial and its derivative are coprime. Over $L=K(r)$,
$r^2=t$, the substitution $y=rY$ identifies $C_L$ with $(C_0)_L$.
The nontrivial inertia element acts on $C_0$ as the hyperelliptic
involution $\iota$. It acts as $-1$ on every prime-to-$3$ Tate module
of the Jacobian: the quotient is $\mathbf P^1$, so
$1+\iota=0$ on its degree-zero Picard variety. Thus $d,c,b$ fail,
and good reduction would contradict $a\Rightarrow d$.

Because the curve is \emph{proper}, base extension from an algebraic
closure of the finite field to $\bar K$ preserves its full geometric
fundamental group. Its full outer $G_K$-action is the constant
$G_{\F_3}$-action twisted by the quadratic character
$\chi:G_K\to\langle\iota\rangle$. For any given splitting $s$,
the character $\chi\circ s$ defines a quadratic twist of $C_0$ over
$\F_3$; ordinary finite Galois descent makes this a smooth proper
genus-two curve. Its full outer action is exactly the restriction along
$s$. This proves $e$ already over $\F_3$. Properness in this argument
is essential; we do not substitute an affine constant curve and assume
invariance of its full wild fundamental group under extension of
algebraically closed fields.
\end{proof}

\begin{claim}{B.24}{Proved}
The same failure of $a\Rightarrow b,c,e$ occurs explicitly at $p=2$.
Over $R=\F_2[[t]]$, let $C$ be the smooth proper normalization of
\[
                    y^2+t z^3y+y=z^9.
\]
Then $C_K$ is a proper genus-four curve with $a,d$ true and full
$b,c,e$ false. Failure of $e$ again holds for every residue splitting
and after every finite residue extension. The ramified character which
detects $c$ here occurs on a cyclic quotient of order four.
\end{claim}
\begin{proof}
Start with the smooth elliptic scheme
\[
 E:\quad y^2+txy+y=x^3.
\]
Its discriminant is $1+t^3$, a unit. Its generic fibre is ordinary
since its Weierstrass coefficient $a_1=t$ is nonzero, whereas its special
fibre $y^2+y=x^3$ is supersingular; the elementary characteristic-two
criterion is \citet[Example 5.26]{BC_MoonenIsogs}.
The divisor of $x$ is $P+Q-2O$, where $P=(0,0)$, $Q=(0,1)$,
and $O$ is the origin at infinity. These are three disjoint sections.
Adjoining $z^3=x$ gives the displayed curve. The normalization is smooth
over $R$: away from the three sections the cover is etale, and at each
section its tame cyclic local normalization, after an etale change
extracting a cube root of the unit, is a smooth parameter with base
parameter its cube. All three ramification indices are three, and
Riemann--Hurwitz gives $2g(C)-2=3(3-1)=6$.

We compute a ramified etale cover of the generic elliptic curve. Put
\[
 x=t^2X+t^{-1},\qquad y=t^3Y+t^{-3},\qquad
 B=t^{-3},\qquad A=t^{-12}+t^{-9}.
\]
Its generic equation becomes
\[
                 Y^2+XY=X^3+BX^2+A.                         \tag{B.24.1}
\]
For a point $(X,Y)$ with $X\ne0$, the tangent slope is
$X+Y/X$. Substitution in the duplication formula gives
\[
                  X(2P)=X(P)^2+A/X(P)^2.
\]
The unique nonzero geometric two-torsion point has $X=0$.
Consequently the points of exact order four have $X^4=A$.
Under twice iterated relative Frobenius their images are the two
primitive points of the etale kernel of
$V^2:E^{(4)}\to E$. Writing their coordinates as $(A,Y_4)$, we have
\[
 Y_4^2+AY_4=A^3+B^4A^2+A^4.
\]
Set $w=Y_4/A+A+B^2+B$. The equation is exactly
\[
                       w^2+w=B=t^{-3}.                       \tag{B.24.2}
\]
The right side has odd pole order three, so this Artin--Schreier
extension is totally ramified of degree two: any root has valuation
$-3/2$. It exchanges the two primitive generators of the etale cyclic
group $\ker V^2(\bar K)\simeq\Z/4$. The relative Frobenius map is
used only to locate this etale kernel; no identification of all
geometric torsion points with points over $K^{\mathrm{sep}}$ is made.
Its pullback along the degree-three map $C\to E$ is geometrically
connected, since three and four are coprime. Thus it gives an
inertia-ramified quotient of $\pi_1(C_{\bar K})^{\mathrm{ab}}$.
This proves failure of $c$ and $b$, despite good reduction.

For completeness, the special curve $y^2+y=z^9$ has $2$-rank zero
directly by Cartier. The different exponent at its unique point over
$z=\infty$ is ten, so $dz$ has divisor $6P_\infty$; a basis of
regular differentials is $dz,z\,dz,z^2dz,z^3dz$. Cartier kills the
even powers and sends $z\,dz$ to $dz$ and $z^3dz$ to $z\,dz$.
Its third iterate is zero. Generically, pullback and norm for the
degree-three map to the ordinary $E$ have composite multiplication by
three, so the etale two-primary Tate module of $E$ is a direct summand
of that of $\Jac(C)$. In particular $f(C_{\bar K})\ge1$.
The special $2$-rank is zero, whereas the generic one is positive.
If a full residue realization existed, prime-to-two specialization,
the properness/genus invariants and Saidi--Tamagawa Theorem 3.9 would
identify it with the special curve, exactly as in B.21. Full mod-two
homology contradicts the two different ranks. This proves the claim.
\end{proof}

The exact identities, Hasse coefficients, smoothness and node checks in
B.21--B.24 are verified in \texttt{code/BC\_equal\_char\_checks.py};
all 35 named checks pass. They verify the explicit examples, whereas
the all-cover specialization argument is proved in B.20.

